%RZ\documentclass[12pt]{article}
\documentclass[letterpaper,titlepage,11pt]{article}

% Added by Roger for preprint colored box
% Highlight text 
%\usepackage{color,soul}
\usepackage{soul}
\usepackage[usenames, dvipsnames]{color}
\usepackage{tcolorbox}
%\newtcolorbox{mybox}{colback=red!5!white,colframe=red!75!black}
\newtcolorbox{mybox}{colback=green!5!white,colframe=green!75!black}

\usepackage{amssymb}
\usepackage{enumitem}

%Added by DH
\usepackage{adjustbox} %Adjust boxed content
\usepackage[colorlinks,urlcolor=blue]{hyperref}
\usepackage{booktabs, multirow}
\usepackage{ragged2e}
\usepackage{array}
\usepackage{amsmath}
\allowdisplaybreaks[1]
\newcolumntype{H}{>{\setbox0=\hbox\bgroup}c<{\egroup}@{}}
\usepackage{algorithm,algpseudocode} 
\usepackage{rotating} % Rotating table
\usepackage{comment}
\usepackage{float}
\usepackage{caption}
\usepackage{subcaption}
\usepackage{changepage,threeparttable} % for wide tables
%\usepackage[numbers]{natbib}
\usepackage[round,semicolon,numbers]{natbib}
\setcitestyle{square,comma,numbers}
\usepackage{enumitem}
%\usepackage{hyperref}
\usepackage{scrextend}
\renewcommand{\algorithmicrequire}{\textbf{Input:}} %redefinecomm
\renewcommand{\algorithmicensure}{\textbf{Output:}}
%RZ\usepackage[affil-it]{authblk}
%RZ\renewcommand*{\Affilfont}{\normalsize}

%RZ
\usepackage{setspace}
%\usepackage[margin=2cm]{geometry}
\usepackage{comment} 
\usepackage{listings,lscape}
\usepackage{lineno} 

%
%       Text dimensions for preprint (Roger). 
%
\setlength{\textwidth}{6.5in}
\setlength{\textheight}{9.0in}
\setlength{\oddsidemargin}{0.0in}
\setlength{\evensidemargin}{0.0in}
\setlength{\topmargin}{0.0in}
\setlength{\headsep}{0.0in}


%New added by DH
\newcommand{\diana}[1]{{\color{purple}#1}}
\usepackage{appendix}
\usepackage{changepage,threeparttable} % for wide tables
\newcommand\colorrow{\rowcolor{blue}\cellcolor{white}}



\title{Enhancing Matheuristics for a Thermal Unit Commitment Problem through Kernel Search and Local Branching}

\author{%
{Uriel I. Lezama-Lope}\\% 
Independent Researcher\\
%Universidad Aut\'onoma de Nuevo Le\'on (UANL)\\ 
%Graduate Program in Electrical Engineering\\%
Cuernavaca, Mexico\\%
E-mail: \textit{urieliram@yahoo.com}\\[1em]
%
%\and
{Roger Z.\ R\'{\i}os-Mercado}\thanks{Corresponding author}\\%
Graduate Program in Electrical Engineering\\%
Universidad Aut\'onoma de Nuevo Le\'on (UANL)\\% 
San Nicol\'as de los Garza, NL, Mexico\\%
E-mail: \textit{roger@optimizer.mx}\\[1em]
%\and
{Diana L. Huerta-Mu{\~n}oz}\\% 
Graduate Program in Systems Engineering\\%
Universidad Aut\'onoma de Nuevo Le\'on (UANL)\\%
San Nicol\'as de los Garza, NL, Mexico\\%
E-mail: \textit{dianahuerta@yalma.fime.uanl.mx}\\
}


\date{June 2025}   

%\setstretch{1.1}
%\renewcommand{\baselinestretch}{1.1}
\begin{document}
\setstretch{1.2}
%\begingroup
%\let\center\justify
%\let\endcenter\justify
\maketitle
%\endgroup

%\linenumbers

%\renewcommand{\baselinestretch}{1.3}
%\setstretch{1.4}
\begin{abstract}
%The Unit Commitment Problem (UCP) consists of the optimal scheduling of electrical generators in a power line system. Typicall, the different variations of the unit commitment problem are hard to solve. 
% In this study, we address a thermal UCP with a staircase cost structure, a formulation that introduces additional complexity. Despite promising results from reformulation strategies, their extended resolution times often make them impractical for real-time applications. To address this issue, we propose a two-phase matheuristic approach. In the first phase, we develop a constructive method to generate high-quality initial solutions. In the second phase, we enhance these solutions using components such as Kernel Search and Local Branching, effectively combining the strengths of mathematical programming and metaheuristics into a hybrid framework. Experimental results in real-world instances demonstrate the effectiveness and value of each of the matheuristic components and that the proposed approach consistently outperforms standard solvers and existing heuristics, particularly for large-scale instances.\\[1em] 

\setstretch{1.3}
%\textcolor{blue}{The Thermal Unit Commitment Problem (UCP) represents a pivotal challenge in the optimization of power system operations, characterized by the need to determine the optimal scheduling of electrical generators to meet fluctuating demand while minimizing production costs. This study addresses the complexities associated with UCP, particularly in the context of staircase cost structures. Despite the encouraging outcomes from the reformulation strategies to the mathematical model, the UCP still presents significant challenges. We propose an innovative two-phase matheuristic framework that first employs a constructive heuristic to generate high-quality initial solutions, followed by an enhancement phase that implements concepts such as Restricted Candidate List (RCL), soft fixing, and optimization strategies utilizing Local Branching and Kernel Search methods. Our approach leverages the synergistic benefits of mathematical programming and metaheuristic strategies, delivering an effective hybrid solution methodology. Empirical evaluations on well-known complex instances from the literature demonstrate that our proposed framework consistently outperforms other state-of-the-art matheuristics as well as off-the-shelf solvers, particularly for medium- and large-scale problem sets. The findings underscore the efficacy of our matheuristic approach, offering significant advancements in the field of UCP resolution and fostering improved reliability in power generation scheduling.} 


The unit commitment problem is an important problem 
arising in the planning and optimization of power system operations.
The problem is characterized by the need to determine the optimal scheduling of electrical generators to meet fluctuating demand while minimizing production costs.
In this study, we address a thermal unit commitment problem under
a staircase cost structure. 
Despite the encouraging outcomes from the reformulation strategies to the mathematical model, this problem still presents significant challenges.
%
Despite the encouraging outcomes from the reformulation strategies to the mathematical model, the UCP still presents significant challenges.
%
We propose a novel two-phase matheuristic framework that first employs a constructive heuristic to generate high-quality initial solutions, followed by an enhancement phase that takes advantage of mechanisms such as local branching and kernel search methods. 
%
%Our approach leverages the synergistic benefits of mathematical programming and metaheuristic strategies, delivering an effective hybrid solution methodology. 
Experimental results on real-world instances demonstrate the effectiveness and value of each of the matheuristic components and that the proposed approach consistently outperforms standard solvers and existing heuristics, particularly for large-scale instances.
%
%Empirical evaluations on well-known complex instances from the literature demonstrate that our proposed method consistently outperforms other state-of-the-art matheuristics as well as off-the-shelf solvers, particularly for medium- and large-scale problem sets. The findings underscore the efficacy of our matheuristic approach, offering significant advancements in the field of UCP resolution and fostering improved reliability in power generation scheduling.

\vspace{1em}
%\noindent\textit{Keywords:} Thermal unit commitment problem; Integer programming; Matheuristics; Kernel search; Local branching.
\noindent\textit{Keywords:} Energy management and planning; Energy system optimization; Thermal unit commitment problem; Optimization algorithms; Matheuristics.
\end{abstract}

%% main text
\section{Introduction}
\label{sec:Introduction}

This research aims to improve the efficiency of solving a thermal unit commitment problem (UCP), a key challenge in power system planning. The UCP involves making decisions about generator operations, such as switching units on or off and determining optimal power output levels, to meet consumer demand while satisfying operational, technical, and economic constraints. As a daily optimization task, it seeks to minimize production costs while adhering to limitations imposed by the electrical system, generator characteristics, and market requirements. Commonly, the UCP is addressed through the formulation of an adequate mixed-integer linear programming (MILP) model.

Despite recent advances on models and algorithms, solving UCPs within limited time frames remains a challenge. To tackle this problem, the purpose of this study is to provide a hybrid heuristic framework that attempts to exploit the advantages of components such as kernel search and local branching with the aim of achieving optimal or near-optimal solutions in the time allowed.

The proposed solution algorithm consists of a two-phase approach, comprising a constructive phase and an improving phase. In the first phase, 
we develop a relax-and-fit heuristic (named HARDUC) that builds feasible solutions to the problem.
Then, in the second phase, we derive a method that aims to improve the solution by adapting local branching~\citep{Fischetti2003} and kernel search ~\citep{Angelelli2010} mechanisms.

The general algorithm and its individual components were fully assessed
in three groups of instances derived from \citet{Kazarlis1996}.
The empirical evaluation included a comparison with
the constructive HGPS heuristic of \citet{Harjunkoski2021} and the
CPLEX solution of the thermal UCP MILP model proposed by \citet{Knueven2020b}.

It was observed that
although the branch-and-bound solver demonstrates strong performance for small-scale problem instances, our proposed method consistently achieves better results for medium- and large-scale instances within the allocated time frame. In particular, our methods reliably generate solutions across all cases, whereas the solver frequently fails to find feasible solutions in medium-, and large-scale instances.

The paper is structured as follows. The most relevant related work
on algorithms, including matheuristics, for the UCP is surveyed in Section~\ref{sec:survey}, \textcolor{blue}{where the specific contribution of the paper is also outlined}.  Section~\ref{sec:Mathematicalmodel}
presents the problem, its assumptions, and the mathematical formulation.
This is followed by Section~\ref{sec:method}
that includes a full and detailed description of the proposed
solution framework and its individual components.
Section~\ref{sec:Experimentalwork}
presents an empirical evaluation of
the algorithm and a comparison with existing work.
Closing remarks, conclusions, and discussion of
future research directions are given in Section~\ref{sec:Conclusions}.



\section{Literature Survey}
\label{sec:survey}

%\color{RedOrange}

%RZ\section{Related work}
%RZ\label{sec:Relatedwork}

%Before continuing, we believe it is necessary to briefly discuss the UCP prior to addressing the matheuristics methods that have been used to solve it.
UCP models are classified into deterministic and stochastic.  In our work, we deal with
a deterministic model, therefore we focus our literature review on deterministic models and approaches. For stochastic UCPs, the reader is referred
to the work by \citet{2019_ijepes_Haberg}, who reviews models and algorithms developed for handling uncertainty in UCPs.

The literature on UCP models is vast \citep{Montero2022}. UCP models have been widely used to manage energy production for many decades. Different variations of UCP models addressing particular situations and assumptions have been studied \citep{Idriss2018}.

UCP models can be classified according to generator operating, electrical network, and system constraints. The first set includes technical constraints related to the generators as power limits, ramps, minimum up and down times, and start-up costs. The second set comprises limits in lines and tie-lines. The last one encompasses meeting demand and load-generation balance. \citet{Anjos2017} outline some examples of these models.

Recent research has established the UCP as a computationally complex challenge. \citet{Anjos2013} affirm that the UCP is an NP-hard, large-scale, non-linear, non-convex, combinatorial optimization problem. On the other hand, \citet{Bendotti2019} assert that the UCP can be transformed into multiple knapsack problems with linking constraints in time and analyze the complexity of the UCP with respect to the number of units and time periods. 
%The UCP is strongly NP-hard, as proven in this study. 


\subsection{Traditional approaches}

In his doctoral thesis, \citet{Morales-Espania2014_thesis} highlights that a significant portion of the UCP literature concentrates on improving the formulation by seeking locally ideal or locally tighter representations for a specific subset of constraints of UCP.  However, the thesis concludes that achieving an excessively tight model becomes useless if computational limitations slow its computability due to the extensive number of variables and constraints involved. \textcolor{blue}{Given the insights provided by these findings, recent research has focused on finding tighter formulations that have shown promising results, particularly in thermal UCP models.  \citet{CarrionArroyo2006} are the pioneers in this modeling approach, proposing a one-binary-variable formulation that significantly reduces the number of variables in the problem. Subsequently, \citet{Ostrowski2012} later proposed a tighter three-binary-variable model and demonstrated that, despite the compactness of the one-binary formulation, the one-binary formulation produces weaker linear relaxations and makes it more difficult to incorporate additional valid inequalities, which ultimately affects computational performance and puts general-purpose solvers at a disadvantage.}

\textcolor{blue}{Subsequently, \citet{Morales-Espania2013} proposed a tight and compact (T\&C) formulation for the thermal UCP model, which aims to minimize operating costs over a specified period by optimizing the commitment and dispatch of generators that use gas, coal, fuel oil, and diesel as primary energy sources. The formulation considers generator capacity limits, minimum up and down times, ramp rates, and variable start-up costs, while ensuring that electricity demand and reserve requirements are met.}

Recent research on UCP models has concentrated on developing more effective Mixed Integer Linear Programming (MILP) formulations that can represent the various components of the problem more tightly and compactly \citep{Morales-Espania2013b,Morales-Espania2016,Yang2017,Yang2021}. These studies have particularly focused on the thermal UCP model. In addition, \citet{Guedes2017} have presented a hydraulic UCP model with T\&C features.

On one hand, a tighter model reduces or narrows down the best solution space to find the feasible solution and helps methods based on branch and bound (B\&B) reach a solution quicker. On the other hand, a more compact model employs fewer constraints and variables and requires fewer computing resources. \citet{Knueven2020b} have compiled the major modern UCP formulations and proposed different models that balance out T\&C features. The new formulations are derived from combining constraints from other models. They also tested the model's performance with instances from electricity markets, showing positive results for tighter models.

These T\&C models have shown promising results because they can be extended with additional constraints as \citet{Nycander2021} that propose a variant based on power and not in energy. This variant of UCP based on power provides a notable advantage by incorporating a more realistic approach to considering generator ramps. It accurately calculates the generation capacity at the beginning and end of each time period. Additionally, it integrates constraints that account for intra-hour reserve requirements, leading to enhanced practicality within the model.

Finally, UCP models have been solved using the branch-and-bound (B\&B) method and its variants \citep{Wolsey2021}. In practice, sophisticated solvers and high computing power are required to solve large-scale problems at reasonable times. Despite this progress, depending on the size and features of specific instances, sometimes the algorithms cannot reach an optimal solution in the time allowed. In this context, metaheuristic methods help find near-optimal solutions. However, some methods based on Genetic Algorithms \citep{Kazarlis1996}, Tabu Search \citep{Mori1996}, and Variable Neighborhoods Search \citep{Todosijevic2016}  have yet to perform well. Nevertheless, recent research has yielded promising results by integrating metaheuristic methods with mathematical programming to address UCPs. Consequently, the next sub-section comprehensively surveys the most recent matheuristic approaches for solving UCP models.

%\color{black}

\subsection{Hybrid heuristic approaches}

A strategy for reducing solution times and improving accuracy in the UCP is the hybrid heuristic methods, known as matheuristics. These strategies cleverly embed exact methods into metaheuristic approaches \citep{2023_itor_Boschetti,Maniezzo2021}. They also take advantage of advances in good MILP formulations \citep{Fischetti2016}. 

The efficiency of branch-and-bound algorithms has been significantly improved by the implementation of heuristic ideas in MILP solvers \citep{Wolsey2021}. Matheuristics can improve solution quality, convergence guarantees, and computational efficiency, leading to better solutions. Constructive heuristics, such as relax-and-fix (R\&F) methods \citep{Wolsey2021}, the feasibility pump (FP) \citep{Fischetti2003}, and the relaxation-induced neighborhood search (RINS) \citep{Danna2005}, use continuous solutions from linear relaxation and rounding to guide the search for feasible integer solutions. This approach improves the quality and efficiency of the solutions by exploring the neighborhoods around relaxed solutions in combinatorial optimization problems.

Other matheuristics inspired by local search algorithms are then used to refine solution quality, such as local branching (LB) strategies introduced by \citet{Fischetti2003}, which define the neighborhoods within the MILP with maximum modifications to the incumbent solution. These strategies iteratively explore neighboring solutions by applying small changes, focusing on improving the incumbent while maintaining feasibility. The approach allows for an effective exploration of the solution space, providing higher-quality solutions by guiding the search toward more optimal regions.

The Kernel Search (KS) matheuristic, introduced by \citet{Angelelli2010}, is inspired by greedy algorithms. The method begins by partitioning the decision variables of an MILP formulation into two sets: a central kernel and multiple buckets. A simplified version of the problem, or sub-MILP, is then solved for each bucket, using the kernel as a starting point. Iteratively, the method adds new variables from the buckets to the kernel, progressively improving the solution. The process continues until all the variables from the buckets have been integrated into the kernel, resulting in an optimal or near-optimal solution with the variables assigned within the kernel.

Following this, a survey of the matheuristic and heuristic methods applied to variations of UCP models is presented. It is worth noting that there is a significant lack of research on the use of matheuristics in the UCP. In Table \ref{tab:MathAuthorsComparison}, the comparison between the works that have addressed the UCP problem with these heuristic approaches is presented. The table lists the main constraints and features of the UCP models addressed in each work and their corresponding solution methods.



\begin{table}[htp]
\caption{Survey of works on UCPs using matheuristics.}\label{tab:MathAuthorsComparison}
\begin{adjustbox}{max width=\textwidth}
\centering
\begin{threeparttable}[b]
\begin{tabular}{>{\RaggedRight}p{3cm}>{\RaggedRight}p{3cm}>{\RaggedRight}p{3cm}>{\RaggedRight}p{1.5cm}>{\RaggedRight}p{3cm}>{\RaggedRight}p{1.5cm}>{\RaggedRight}p{1.5cm}>{\RaggedRight}p{2cm}>{\RaggedRight}p{2cm}>{\RaggedRight}p{1.5cm}>{\RaggedRight}p{1.5cm}}\toprule
{Article} & {Method} &{UCP type} & {PC} & {SUC} & {SDC} & {RR} & {GL} & {RUD} & {SR} & {MUDT} \\\cmidrule{1-11}
\citet{Fayzur2014} 2014 &LB-ILA, PSO-ILA &thermal &Q &Cold and hot start-up & \checkmark & \checkmark & \checkmark & \checkmark &  & \checkmark \\\cmidrule{1-11}
\citet{Todosijevic2016} 2016 &VNS-linear programming &thermal &Q &Cold and hot start-up & \checkmark & \checkmark & \checkmark &  &  & \checkmark \\\cmidrule{1-11}
\citet{Saboia2016} 2016 &LB &stochastic thermal compact formulation &L &Exponential start-up cost & \checkmark &  & \checkmark & \checkmark &  & \checkmark \\\cmidrule{1-11}
\citet{Dupin2016} 2016 &R\&F, VNS with MIP neighborhoods. LB, RINS. &discrete thermal, real-time &L &Fixed cost &  & \checkmark & \checkmark &  &  & \checkmark \\\cmidrule{1-11}
\citet{Dupin2018} 2018 &R\&F-B\&B &discrete thermal, real-time &L &Fixed cost &  & \checkmark & \checkmark & \checkmark &  & \checkmark \\\cmidrule{1-11}
\citet{Santos2020} 2020 &FP, LB &Hydrothermal &L &Fixed cost & \checkmark &  & \checkmark & \checkmark & \checkmark & \checkmark \\\cmidrule{1-11}
\citet{Harjunkoski2021} [17] 2021 &R\&F &thermal &L &Fixed cost & \checkmark & \checkmark & \checkmark & \checkmark &  & \checkmark \\\cmidrule{1-11}
\citet{Liu2023} 2023 &Learning-assisted MIP & stochastic thermal & L & Variable start-up cost &  &  & \checkmark & \checkmark & \checkmark &  \checkmark\\\cmidrule{1-11}
{This work} &R\&F, LB, KS. &thermal T\&C model &S &Variable start-up cost & \checkmark & \checkmark & \checkmark & \checkmark & \checkmark & \checkmark \\
\bottomrule
\end{tabular}
\begin{tablenotes}
    \item PC: Production cost; SUC: Start-up cost; SDC: Shut-down cost; RR: Reserve requirements; GL: Gen limits; RUD: Ramp up/down; SR: Start-up/shut-down ramps; MUDT: min up/down times.
	\item Q: quadratic; L: linear; S: staircase.
\end{tablenotes}
\end{threeparttable}
\end{adjustbox}
\end{table}


%\citet{Xavier2021} 2021 &Learning-assisted MIP &SCUC &Q &Fixed and variable start-up cost & \checkmark &  & \checkmark & \checkmark & \checkmark & \checkmark \\\cmidrule{1-11}

First, \citet{Fayzur2014} present two matheuristic approaches to solve a thermal UCP. The first approach combines the original version of LB from \citet{Fischetti2003} and an iterative linear approximation (ILA) method. The second approach combines particle swarm optimization (PSO) and ILA, incorporating a solver in an iterative process.
The UCP problem addressed in the study initially involves quadratic production costs, which are later linearized with ILA and solved using the LB method. The model's constraints include power balance, spinning reserve requirements, minimum up and down times, production limits, and ramps \citep{Viana2013}.
The study solves previously unsolvable instances and achieves faster optimal results for larger instances. The proposed LB-ILA algorithms significantly reduce the average CPU time for larger instances by 40\%-56\%. However, these algorithms perform worse for smaller instances due to unnecessary exploration of the neighborhoods.

\citet{Saboia2016} solve a stochastic thermal network-constraint UPC using the LB approach combined with an iterative approach that considers transmission lines flow limits. 
They dynamically introduce violated flow limit constraints, embedding them into the nodes of the LB scheme. This iterative procedure involves fixing the optimal integer solution achieved at each node. The DC power flow is then evaluated iteratively, continuously checking for violations of line capacity. The process continues until no further line capacity violations are observed. This model is a new compact MILP formulation with basic constraints for thermal UCP. The method is tested on one instance.

\citet{Todosijevic2016} propose a hybrid approach that combines variable neighborhood search (VNS) with mathematical programming to address a UCP. 
In their method, the commitment of generators, i.e., deciding which ones will be turned on and off, is determined using VNS. Additionally, they solve an economic dispatch problem (calculus of each generator's power level) for each period by formulating the dispatch as a linear programming problem.
The results obtained by the proposed approach outperformed considerable metaheuristics reported in the literature for solving the UCP. Furthermore, the solutions achieved were very close to the optimal solutions obtained with the CPLEX solver. However, the paper lacks information regarding the CPU times of the solver, making it impossible to compare the computational efficiency of VNS directly.
Nevertheless, it is important to note that the mathematical model used in the study considers only constraints of demand, reserves, power limits, minimum up and down times, and hot and cold start-up costs. However, the model does not incorporate ramps, which are essential constraints in real-life UCP usage. 
Furthermore, a particular feature of their UCP is that the production cost of the generators is quadratic.

\citet{Dupin2016} use a matheuristic method to solve a discrete thermal UCP from programming thermal generators in the real-time electricity French system. 
First, in a constructive phase, they use R\&F strategies to find one initially good solution. 
In the subsequent improvement phase, a VNS algorithm conducts a local search by iteratively exploring B\&B solutions within neighborhoods defined within the solution space confined by the MILP model. In other words, these neighborhoods are defined within the MILP using various heuristics suggested by the authors, such as RINS or LB strategies.

The mathematical model used by \citet{Dupin2017} comprised of constraints such as demand, reserves, minimum up and down times \citep{Rajan2005}, and fixed start-up costs. 
Note that in the discrete thermal UCP, the limits of generation constraints are not applicable. Instead, discrete power levels and transition constraints exist between them.
However, the model does not incorporate ramps.
The method aims to find the best solutions for the discrete thermal UCP on a 15-minute time limit from its application in a real-time scenario.
The results of the metaheuristic method are better than the MILP solution, which was obtained by the solver in the allotted time frame. 

In follow-up work, \citet{Dupin2018} improved the discrete thermal UCP formulation with tight and compact (T\&C) features and extended the model to include min-stop ramping for a discrete thermal UCP.  To address the problem, the authors propose several constructive matheuristics designed to generate good-quality feasible solutions. These solutions can be effective initial solutions for a branch-and-bound (B\&B) algorithm. Among all the ingenious constructive methods presented, they implemented one with the KS ideas proposed by \citet{Guastaroba2017}, i.e., the segregation of generators based on their marginal costs, eliminating units with higher production costs. Another variant of their heuristic prioritizes utilizing low-cost units as the main production basis, while the remaining units constitute the kernel of the KS algorithm. However, the authors neither utilize the complete KS method nor provide specific details about the implementation
of the method. The proposed method produces precise, high-quality solutions, outperforming the B\&B method applied to the MILP model in the allotted time limit.
The authors tested the method with more than 600 instances from the French electrical system.


According to the investigation undertaken by the \citet{Santos2020}, who present the model and method used to plan operations for energy generation in Brazil, their work considers thermal generation and hydro generation and limitations on the electrical network. They used a variant of the FP from \citet{Fischetti2005} adapted for binary variable decisions. They also created an objective function adapted with a mathematical term that calculates the Hamming distance. The researchers opted for a heuristic method due to the limitations in time to solve the UCP, which should be solved daily for a time horizon of seven days. Considering the dimensions of their problem, an exact method would not yield a timely resolution within two hours. The authors report using the same version of LB as that of \citet{Saboia2016}.

\textcolor{blue}{
\citet{Liu2023} address a thermal stochastic UCP with network constraints in thermal systems with high renewable penetration using a straightforward learning method that does not rely on complex machine learning techniques. They analyze historical daily commitment patterns to fix stable binary variables and eliminate rarely used ramping constraints. This simple, rule-based approach effectively reduces the problem size and computational time without sacrificing solution quality, providing a practical alternative to more complex learning-based methods.
\citet{Xavier2021} propose a learning-based approach for  security UCP that predicts redundant network constraints, generates warm-start solutions, and identifies affine subspaces to reduce problem size and accelerate computation.
Although both \citet{Liu2023} and \citet{Xavier2021} leverage learning from multiple daily real instances, \citet{Liu2023} fix binary commitment variables and removes underused ramping constraints, whereas \citet{Xavier2021} emphasize predicting redundant network constraints, generating warm-starts, and identifying affine subspaces to reduce problem size and computation time.
}

In the study conducted by \citet{Harjunkoski2021}, an R\&F heuristic for the UCP is devised. The proposed method starts with solving a linear relaxation of the problem. Subsequently, they analyzed the product of the commitment binary variable with each generator's power level. The commitment binary variables were fixed and solved as an MILP problem if the product exceeded the minimum feasible generation level. This constructive method provides a warm start to the B\&B algorithm and accelerates the UCP solution process.



%RZ In this section, we discuss the key distinctions between our work and the related studies, emphasizing the scientific contributions of our research.

\subsection{Positioning our work in the literature}

Regarding the mathematical modeling of the problem, our UCP is a deterministic thermal model with a staircase cost function. 
Furthermore, our model incorporates valid inequalities in the cost constraints, making the formulation tighter.
This cost function assigns different price levels for each generation level interval, 
aligning the model more closely with the cost modeling of electricity markets. 
Other works address similar problems with linear costs \citet{Saboia2016,Dupin2016,Dupin2018,Harjunkoski2021} and quadratic costs \citet{Fayzur2014,Todosijevic2016}.
Another notable difference is that the modeling of generator startup costs is simplified in most of the papers to a fixed startup cost \citet{Dupin2016,Dupin2018,Santos2020,Harjunkoski2021} 
or hot and cold starts \citet{Fayzur2014,Todosijevic2016}, or, at best, exponential startup costs depending on the time the generator has been off \citet{Saboia2016}. 
In our case, we use variable startup cost constraints \citet{Morales-Espania2013} with tighter features.

The main difference with the other works is that, although they also solve a thermal UCP, they address different variants. 
For example, \citet{Dupin2016,Dupin2018} tackle a discrete thermal UCP for a real-time horizon with limited forward scheduling periods. 
\cite{Santos2020} deal with a hydrothermal UCP problem, while \citet{Saboia2016} solve a stochastic UCP thermal problem. 
Although \citet{Fayzur2014} present a hybrid LB method with cost linearization, the version they report is the original one developed by \citet{Fischetti2003}. 
Alternatively, \citet{Todosijevic2016} apply VNS to solve the problem and incorporates linear programming to address a subproblem. 
Similarly, \citet{Dupin2016} employ another VNS, and some of the neighborhoods they define incorporate the concepts of LB and RINS matheuristics. 
Moreover, it is worth noting that only the local search definition from the LB method is utilized without involving the branching phase. 
Additionally, \citet{Saboia2016} modify the original LB to include the power network.  
%
In their research, \citet{Santos2020} report using a variant of FP as a constructive strategy and LB as an improvement strategy. 
Furthermore, the authors state that they use the same version of LB as \citet{Saboia2016}.

\textcolor{blue}{
Similar to matheuristic strategies that fix variables with persistent patterns and remove underused ramping constraints, \citet{Liu2023, Xavier2021}  simplify the MILP problem using information extracted from past instances. In contrast to our approach, which does not rely on learning from previously solved instances, these learning-based works show how historical solutions can guide effective problem-size reduction in stochastic and security-constrained UCP settings.
}

The first paper to report on the use of the KS for UCP is \citet{Dupin2018},
where they present their findings on incorporating the KS concept as a variable-fixing criterion into a variable fixing strategy. 
This strategy introduces methods to select specific variables to fix a priori to obtain a reduced problem that the MILP solver can subsequently solve. 
However, it is important to emphasize that their work lacks a complete implementation of the KS method.

Furthermore, \citet{Harjunkoski2021} propose a constructive method based on R\&F, where the initial feasible solution provides a warm start to the B\&B algorithm, resulting in an accelerated UCP solution process. 

\subsection*{Summary of research contributions}

In terms of solution methodologies, we proposed a novel constructive method that aims to compete with the approach of \citet{Harjunkoski2021} and the solver. The solution calculated by the constructive method provides the first solution to the improving method.

In the improvement phase, we have developed a unique version of KS; our version introduces several significant innovations compared to the original method proposed by \citet{Angelelli2010}. First, our KS eliminates the need for kernel construction, as a preliminary constructive stage already provides the kernel. Second, we employ the statistical rule of Sturges to determine the number of buckets in KS, a refinement that optimizes the method's performance. Third, our focus is on the dominant variables of the problem, which enhances computational efficiency. Fifth, our method calculates the reduced costs by fixing kernel variables while leaving the remaining variables free, using the linear relaxation of the problem. Finally, our kernel expansion process continues even after the buckets have been processed. If time permits, the kernel is reconfigured using the latest solution, and the expansion process is restarted, allowing for further refinement of the results.

%RZ\hl{RZ No se menciona nada aqui sobre local branching??}
\textcolor{black}{
Furthermore, we introduce four specialized versions of the Local Branching (LB) method, each incorporating refinements to enhance performance. Our adaptations present significant differences compared to the original approach by \citet{Fischetti2003}. The first procedure, LB1, defines a restricted candidate list (RCL) based on variables that are nonzero in the constructive method (HARDUC) but fall below the minimum generator power threshold. Additionally, LB1 employs a soft-fixing constraint that retains at least 90\% of the binary support variables in the solution. The second procedure, LB2, removes the soft-fixing constraint while maintaining the RCL based on the HARDUC rule. The third procedure, LB3 adheres closely to the original LB method but does not use an RCL or soft-fixing constraints, limiting its local search to binary support variables and their complements. Finally, LB4 selects the RCL based on negative reduced costs, improving the efficiency of the search process. A key feature of our LB adaptations is that they operate solely on the dominant variable $u_{g,t}$ to define the binary support, unlike the generic LB approach, which considers all binary variables.
}

\section{Description of the Thermal Unit Commitment Problem}
\label{sec:Mathematicalmodel}


\subsection{Problem definition and assumptions}


The Unit Commitment Problem (UCP) aims to minimize the total operational cost of a set of generators $g \in \mathcal{G}$ over a planning horizon $t \in \mathcal{T}$. The objective function considers several cost components: production costs $c_{gt}^{\text{p}}$, fixed costs $C_g^{\text{R}}$, start-up costs $c_{gt}^{\text{SU}}$, and shutdown costs $c_{gt}^{\text{SD}}$. Achieving this minimization requires adherence to constraints that reflect the physical and operational limits of the system, ensuring feasible and optimized generator scheduling. These constraints govern aspects such as power generation bounds, minimum uptime and downtime requirements, ramping limits, demand-reserve satisfaction, and staircase cost modeling, and are represented using binary and continuous variables for generator states and power outputs.

The constraints of the UCP encapsulate the physical and operational characteristics of power generation systems, ensuring solutions are both feasible and realistic. Each constraint represents a specific aspect of generator behavior or grid requirements:

\begin{description}
    \item[Power limits:] Each generator must operate within its designed capacity, defined by the bounds $\underline{P}_g u_{gt} \leq p_{gt} \leq \overline{P}_g u_{gt}$ for all $g \in \mathcal{G}$ and $t \in \mathcal{T}$. This prevents overloading and maintains efficient operation.

\item[Minimum uptime and downtime:] Generators must stay online for at least $UT_g$ periods after starting up and offline for at least $DT_g$ periods after shutting down. These constraints, $\sum_{\tau=t}^{t+UT_g-1} u_{g\tau} \geq UT_g v_{gt}$ and $\sum_{\tau=t}^{t+DT_g-1} (1 - u_{g\tau}) \geq DT_g w_{gt}$, prevent excessive wear and ensure reliable operation.

\item[Ramping constraints:] The rate at which power output changes is limited by ramp-up and ramp-down constraints, defined as $p_{gt} - p_{g(t-1)} \leq RU_g$ and $p_{g(t-1)} - p_{gt} \leq RD_g$. These capture the physical limitations of generator components.

\item[Start-up and shutdown costs:] These costs, $c_{gt}^{\text{SU}}$ and $c_{gt}^{\text{SD}}$, depend on generator downtime and reflect additional fuel use and equipment wear during transitions.

\item[Demand and reserve requirements:] Total generation must satisfy demand $D_t$ and maintain a reserve margin $R_t$, expressed as $\sum_{g \in \mathcal{G}} p_{gt} \geq D_t + R_t$ for all $t \in \mathcal{T}$. This ensures grid stability and reliability.

\item[Staircase cost structure:] Production costs vary staircase with output, modeled as $p_{gt} = \sum_{l \in \mathcal{L}_g} p_{gt}^l$, where $p_{gt}^l$ represents the power produced in each segment $l$. This mirrors fuel efficiency dynamics and captures operational states.

\item[Initial conditions:] The initial states and outputs of generators are predefined, with $u_{g0}, p_{g0}, v_{g0}, w_{g0}$ specified to ensure a seamless transition from previous operations.

\item[System constraints:] These maintain the balance between generation, demand, and reserves, ensuring grid-wide operational reliability.
\end{description}

\subsection{Mixed-integer linear programming model
}
This section introduces a comprehensive Thermal Unit Commitment (TUCP) formulation that is a robust benchmark for evaluating subsequent matheuristic methods. Our model incorporates diverse constraints, including power limits, minimum uptime/downtime, ramp capabilities, variable start-up costs, and demand-reserve requirements, while also narrowing down solution space with staircase cost production and valid inequalities. 

\textcolor{blue}{The TUCP formulation used here is deterministic and excludes network and contingency constraints. While standard UCP models are usually tractable for off-the-shelf solvers, prior studies \citep{CarrionArroyo2006,Morales-Espania2013,Ostrowski2012} show that synthetic instances with high structural symmetry can severely hinder branch-and-bound performance. We therefore use such instances to rigorously assess both the solver and the proposed matheuristics, allowing us to test whether hybridizing heuristics with an exact solver yields performance gains over the solver alone. For now, our focus is on validating the effectiveness of the matheuristics methods for the core UCP before moving to more complex, security-constrained formulation.}


\subsubsection*{Notation}
A summary of sets, indices, parameters, and decision variables is enlisted for quick reference.

\vspace{1em}
\noindent\textit{Sets and indices of the power system}:%\vspace{-1.0em}
\begin{labeling}{$\mathcal{G}^{\text{TEX}}$}
\setlength\itemsep{0.1em}
\item[$\mathcal{G}$] Set of generators ($g\in\mathcal{G}$)
\item[$\mathcal{G}^{1}$]  Subset of generators ($g\in\mathcal{G}$) ; (if $UT_{g}=1$)
\item[$\mathcal{G}^{1*}$] Subset of generators ($g\in\mathcal{G}$) ; (if $UT_{g}=1 \ \text{and} \ \textit{SU}_{g} \neq \textit{SD}_{g} $)
\item[$\mathcal{G}^{>1}$] Subset of generators ($g\in\mathcal{G}$) ; (if $UT_{g}>1$)
\item[$\mathcal{T}$] Set of time periods in the planning horizon;  ($t\in \mathcal{T}$)
\item[$\mathcal{T}^0$] $=\mathcal{T} \setminus \{ |\mathcal{T}| \}$ Set of time periods in the planning horizon except the last one
\end{labeling}

\vspace{1em}
\noindent\textit{Sets and indices of cost-related aspects}:%\vspace{-1.0em}
\begin{labeling}{$\mathcal{B}^{\text{ABCD}}$}
\setlength\itemsep{0.1em}
\item[$\mathcal{S}_g$] Set of start-up cost curve segments for a generator ($g\in\mathcal{G}$) from hottest ($s=1$) to coldest ($s=|\mathcal{S}_g|$); ($s \in \mathcal{S}_g$).  
\item[$\mathcal{L}_g$] Stairwise production cost intervals for generator ($g\in\mathcal{G}$); ($l\in \mathcal{L}_g$) 
\end{labeling}

\vspace{1em}
\noindent\textit{Parameters}:%\vspace{-1.0em}
\begin{labeling}{$a_{1,g},a_{2,g},a_{3,g}$}
\setlength\itemsep{0.1em}
\item[$C^{\text{R}}_{g}$] Minimum operating cost of generator ($g\in\mathcal{G}$) that works at least at minimum power $\underline{P}_{g}$; in \$. It comprises the no-load fixed cost plus the cost of producing at $\underline{P}_{g}$, i.e., $C^{\text{R}}_{g}= C^{\text{fix}}_{g} + C^{1}_{g}\underline{P}_{g}$
\item[$C^{l}_{g}$] Cost coefficient for stairwise segment ($l\in\mathcal{L}_g$) of generator ($g\in\mathcal{G}$) that works at least at minimum power $\underline{P}_{g}$; in \$/MWh
\item[$C^{\text{S}}_{g,s}$] Start-up cost for generator ($g\in\mathcal{G}$) with a set of segment time ($s\in\mathcal{S}_{g}$); it determines the starting cost by locating a cost in a segment time $s$ in the intervals $[\underline{T}_{g,s},\overline{T}_{g,s})$; in \$/h
\item[$C^{\text{SD}}_{g}$] Shut-down cost for generator ($g\in\mathcal{G}$); in \$ (equal to zero in all test instances)
\item[$\textit{De}_{t}$] Energy load demand in period ($t\in\mathcal{T}$); in MWh %Nycander
\item[${p}_{g,0}$] Power output of a generation ($g\in\mathcal{G}$) at time 0; in MW
\item[$\overline{P}_{g}$,$\underline{P}_{g}$] Maximum and minimum generation value of generator ($g\in\mathcal{G}$); in MW
\item[$\overline{P}_{g}^{l}$] Maximum power available for staircase segment ($l\in\mathcal{L}_g$) of generator ($g\in\mathcal{G}$); in MW
\item[$R_{t}$] System-wide spinning reserve requirement in period ($t\in\mathcal{T}$); in MW
\item[$\textit{RD}_{g}$] Ramp-down rate is the capacity of generator ($g\in\mathcal{G}$) to decrease power between two consecutive periods; in MW/h
\item[$\textit{RU}_{g}$] Ramp-up rate of generator ($g\in\mathcal{G}$) to increase power between two consecutive periods; in MW/h
\item[$|\mathcal{S}_g|$] Number of segments in the set $\mathcal{S}_g$ 
\item[$\textit{SU}_{g}$] Start-up rate for a generator ($g\in\mathcal{G}$); in MW/h
\item[$\textit{SD}_{g}$] Shut-down rate for a generator ($g\in\mathcal{G}$); in MW/h
\item[$T^{\text{RU}}_g$] Time that a generator ($g\in\mathcal{G}$) spends ramping to go from $\textit{SU}_{g}$ to $\overline{P}_{gt}$; in h
\item[$T^{\text{RD}}_g$] Time that a generator ($g\in\mathcal{G}$) spends ramping to go from $\overline{P}_{gt}$ to $SD_{g}$; in h
\item[$\textit{TC}_{g}$] Time offline after a generator ($g\in\mathcal{G}$) turned into cold
\item[$\underline{T}_{g,s}$] Start of start-up cost segment ($s\in\mathcal{S}_g$), respectively; in h, (i.e. $\underline{T}_{g,1}= DT_{g}, \underline{T}_{g,S_{g}}=\textit{TC}_{g}$
\item[$\textit{UT}_{g}$,$\textit{DT}_{g}$] Minimum up/down time for a generator ($g\in\mathcal{G}$); in h
\item[${u}_{g,0}$] Status of generator ($g\in\mathcal{G}$) at time 0
\item[$U_{g}$] Number of periods generator ($g\in\mathcal{G}$) is required to be online at $t=1$ ; in h
\item[$D_{g}$] Number of periods generator ($g\in\mathcal{G}$) is required to be offline at $t=1$ ; in h
\end{labeling}

\vspace{1em}
\noindent\textit{Binary variables}:%\vspace{-1.0em}
\begin{labeling}{XXXX}
\setlength\itemsep{0.1em}
\item[${u}_{gt}$] Equal to 1 if generator ($g\in\mathcal{G}$) is online in period ($t\in\mathcal{T}$), and 0 otherwise 
\item[$v_{gt}$] Equal to 1 if generator ($g\in\mathcal{G}$) starts up at the beginning of period ($t\in\mathcal{T}$), and 0 otherwise
\item[$w_{gt}$] Equal to 1 if generator ($g\in\mathcal{G}$) is shut-down at the beginning of period ($t\in\mathcal{T}$), and 0 otherwise
\item[$\delta_{gts}$] Equal to 1 if generator ($g\in\mathcal{G}$) have a start-up type $s \in \mathcal{S}$ in period ($t\in\mathcal{T}$), and 0 otherwise 
\end{labeling}

\vspace{1em}
\noindent\textit{Real variables}:%\vspace{-1.0em}
\begin{labeling}{$Loss^{MP}_{t}$}
\setlength\itemsep{0.1em}
\item[$c_{gt}^{\text{p}}$] Production cost over $\underline{P}$ of generator ($g\in\mathcal{G}$)  in period ($t\in\mathcal{T}$); in \$
\item[$c_{gt}^{\text{SD}}$] Shut-down cost of generator ($g\in\mathcal{G}$) in period ($t\in\mathcal{T}$); in \$
\item[$c_{gt}^{\text{SU}}$] Start-up cost of generator ($g\in\mathcal{G}$) in period ($t\in\mathcal{T}$); in \$
\item[$p_{gt}$] Amount of power a generator ($g\in\mathcal{G}$) produces in period ($t\in\mathcal{T}$), in MW 
\item[$p'_{gt}$] Amount of power above minimum $\underline{P}_g$ that a generator ($g\in\mathcal{G}$) produces in period ($t\in\mathcal{T}$), in MW 
\item[$\overline{p}_{gt}$] Maximum power available from generator ($g\in\mathcal{G}$) produces in period ($t\in\mathcal{T}$), in MW 
\item[$\overline{p}'_{gt}$] Maximum power available above minimum from generator ($g\in\mathcal{G}$) produces in period ($t\in\mathcal{T}$), in MW 
\item[${p}^{l}_{gt}$] Power from staircase segment ($l\in\mathcal{L}_g$) 
 from generator ($g\in\mathcal{G}$) produces in period ($t\in\mathcal{T}$), in MW 
\item[$r_{gt}$] Spinning reserves provided by a generator ($g\in\mathcal{G}$) in period ($t\in\mathcal{T}$), in MW 
\end{labeling}

The parameters $C_{gl}^{v}$ and $C_{gl}^{w}$  are calculated as follows \citep{Knueven2020b}. 
\begin{equation}
    C_{gl}^{v} =
    \left\{
     \begin{array}{ll}
       0                                           & \text{if  } \overline{P}_{g}^{l} \leq \textit{SU}_{g} \text{}\\
       \overline{P}_{g}^{l}-\textit{SU}_{g}                 & \text{if  } \overline{P}_{g}^{l-1}<\textit{SU}_{g}<\overline{P}_{g}^{l}\text{}\\
       \overline{P}_{g}^{l}-\overline{P}_{g}^{l-1} & \text{if  } \overline{P}_{g}^{l-1} \geq \textit{SU}_{g}  \text{}\\
     \end{array}
    \right.   \nonumber
\end{equation}
\begin{equation}
  C_{gl}^{w} =
    \left\{
     \begin{array}{ll}
       0                                           & \text{if  } \overline{P}_{g}^{l} \leq \textit{SD}_{g} \text{}\\
       \overline{P}_{g}^{l}-\textit{SD}_{g}                 & \text{if  } \overline{P}_{g}^{l-1}<\textit{SD}_{g}<\overline{P}_{g}^{l}\text{}\\
       \overline{P}_{g}^{l}-\overline{P}_{g}^{l-1} & \text{if  } \overline{P}_{g}^{l-1} \geq \textit{SD}_{g} \text{} \\
     \end{array}
    \right.  \nonumber
\end{equation}


Finally, $T^{\text{RU}}_{g}$ and $T^{\text{RD}}_{g}$ are also computed as follows \citep{Knueven2020b}:
%Finally, $T^{\text{RU}}_{g}$ is also computed as follows \citep{Knueven2020b}.
\begin{equation}
T^{\text{RU}}_{g}=\left\lfloor \frac{\overline{P}_{g}-\textit{SU}_{g}}{\textit{RU}_{g}}  \right\rfloor, \nonumber
\end{equation}
\begin{equation}
T^{\text{RD}}_{g}=\left\lfloor \frac{\overline{P}_{g}-\textit{SD}_{g}}{\textit{RD}_{g}}  \right\rfloor. \nonumber
\end{equation}


\subsubsection*{Mathematical formulation}
The MILP model is given by:
\begin{alignat}{3}
%& \text{Objective function:}  \nonumber
& \text{(TUCP) Minimize}\quad \sum\limits_{t\in\mathcal{T}} \sum\limits_{g\in \mathcal{G}} C_{g}^{\text{R}}u_{gt} + c_{gt}^{\text{p}} +  c_{gt}^{\text{SU}} + c_{gt}^{\text{SD}} \qquad &&  \label{eq:obj}
\\
& \text{Subject to: }  && \nonumber \\
&\sum_{i=t-UT_{g}+1}^{t} v_{gi} \leq  u_{gt}
&&  g\in\mathcal{G}, \, t\in \{ UT_g, \ldots, |\mathcal{T}| \},  \label{MUT}\\
%
& \sum_{i=t-DT_{g}+1}^{t} w_{gi}\leq 1-u_{gt}
&& g\in\mathcal{G}, \, t\in \{ DT_{g},...,|\mathcal{T}| \}, \label{MDT}\\
%
& \sum_{i=1}^{\min\{ U_{g},|\mathcal{T}| \}} u_{gi} = \min \{ U_{g},|\mathcal{T}| \}
&&  g\in\mathcal{G} \label{MUT2}\\
%
& \sum_{i=1}^{\min\{ D_{g},|\mathcal{T}| \}} u_{gi} = 0
&& g\in\mathcal{G} \label{MDT2}\\
%
& u_{gt}-u_{g,t-1} = v_{gt}-w_{gt}
&& g\in\mathcal{G}, \, t\in\mathcal{T} \label{garver}\\
%
& p'_{gt}+r_{gt}\leq \left( \overline{P}_{g} - \underline{P}_{g} \right)u_{gt} - \left( \overline{P}_{g}-\textit{SU}_{g} \right) v_{gt} && \nonumber\\ 
& \quad -\left[ \textit{SU}_{g} - SD_{g} \right]^{+} w_{g,t+1} && g\in\mathcal{G}^{1}, \, t\in \mathcal{T}^0\label{genlim1}\\
%
& p'_{gt}+r_{gt}\leq \left( \overline{P}_{g} - \underline{P}_{g} \right)u_{gt} - \left( \overline{P}_{g}-\textit{SD}_{g} \right) w_{g,t+1} && \nonumber\\
& \quad -\left[ \textit{SD}_{g} - \textit{SU}_{g} \right]^{+} v_{gt}
&& g\in \mathcal{G}^{1}, \, t \in \mathcal{T}^0 \label{genlim2}\\
%
& \overline{p}_{gt} \leq \overline{P}_{g}u_{gt} - \left(  \overline{P}_{g}-\textit{SD}_{g} \right) w_{g,t+1} && \nonumber\\ 
& \quad - \sum_{i=0}^{\min\{ \textit{UT}_g-2, \, T^{\text{RU}}_g \}} \left(  \overline{P}_{g} -\textit{SU}_g-i\,\textit{RU}_g\right)v_{g,t-i}
&& g\in\mathcal{G}, \, t\in \mathcal{T}^0 \label{genlim3}\\
%
& p'_{gt}+r_{gt}\leq \left( \overline{P}_{g} - \underline{P}_{g} \right) u_{gt} - \left( \overline{P}_{g}-\textit{SU}_{g} \right) v_{gt} && \nonumber\\ 
& \quad - \left( \overline{P}_{g} - \textit{SD}_{g} \right) w_{g,t+1}
&& g\in\mathcal{G}^{>1}, \, t\in \mathcal{T}^0 \label{startramp}\\
%
& p'_{gt}+r_{gt}\leq \left( \overline{P}_{g} - \underline{P}_{g} \right) u_{gt} - \left( \overline{P}_{g}-\textit{SU}_{g} \right) v_{gt}
&& g\in\mathcal{G}^{1}, t\in\mathcal{T} \label{shutdownramp}\\
%
& p'_{gt}+r_{gt}\leq \left( \overline{P}_{g} - \underline{P}_{g} \right) u_{gt} - \left( \overline{P}_{g}-\textit{SD}_{g} \right) w_{g,t+1}
&& g\in\mathcal{G}^{1}, t \in \mathcal{T}^0 \label{shutdownramp2}\\
%
& \overline{p}'_{gt}-p'_{g,t-1}\leq (\textit{SU}_{g}-\underline{P}_{g} -RU_{g})v_{gt}+RU_{g}u_{gt} 
&& g\in\mathcal{G}, t\in\mathcal{T} \label{rampup}\\
%
& p'_{g,t-1}-p'_{gt}\leq (\textit{SD}_{g}-\underline{P}_{g}-RD_{g})w_{gt}+RD_{g}u_{g,t-1} \quad
&& g\in\mathcal{G},t\in\mathcal{T} \label{rampdown}\\
%
& p^{l}_{gt}\leq \left( \overline{P}^{l}_{g} - \overline{P}^{l-1}_{g} \right) u_{gt}
&& g\in\mathcal{G}, t\in\mathcal{T}, l\in\mathcal{L}_g \label{Staircase1}\\
%
& \sum_{l\in \mathcal{L}_g} p^{l}_{gt}=p'_{gt}
&& g\in\mathcal{G},t\in\mathcal{T} \label{Staircase2}\\
%
& \sum_{l\in \mathcal{L}_g} C^{l}_{g}p^{l}_{gt}=c_{gt}^{\text{p}}
&& g\in\mathcal{G},t\in\mathcal{T} \label{Staircase3}\\
%
& p'_{gt}\leq \left( \overline{P}_{g} - \underline{P}_{g} \right) u_{gt}
&& g\in\mathcal{G},t\in\mathcal{T} \label{Staircase4}\\
%
& p^{l}_{gt}\leq \left( \overline{P}^{l}_{g} - \overline{P}^{l-1}_{g} \right) u_{gt} - C^{v}_{gl}v_{gt} -C^{w}_{gl}w_{g,t+1} &&  g\in\mathcal{G}^{>1}, \, t\in \mathcal{T}^0, \, l\in\mathcal{L}_g \label{Staircase5}\\
%
& p^{l}_{gt}\leq \left( \overline{P}^{l}_{g} - \overline{P}^{l-1}_{g} \right) u_{gt} - C^{v}_{gl}v_{gt}
&& g\in\mathcal{G}^{1},t\in\mathcal{T},l\in\mathcal{L}_g \label{Staircase6}\\
%
& p^{l}_{gt}\leq \left( \overline{P}^{l}_{g} - \overline{P}^{l-1}_{g} \right) u_{gt} -C^{w}_{gl}w_{g,t+1} 
&&  g\in\mathcal{G}^{1},  \, t\in \mathcal{T}^0, \, l\in\mathcal{L}_g \label{Staircase7}\\
%
& p^{l}_{gt}\leq \left( \overline{P}^{l}_{g} - \overline{P}^{l-1}_{g} \right) u_{gt} - C^{v}_{gl}v_{gt} && \nonumber\\
& \quad - \left[ C^{v}_{gl} -C^{w}_{gl} \right]^{+} w_{g,t+1} 
&& g\in\mathcal{G}^{1*}, \, t\in \mathcal{T}^0, \, l\in\mathcal{L}_g \label{Staircase8}\\
%
& p^{l}_{gt}\leq \left( \overline{P}^{l}_{g} - \overline{P}^{l-1}_{g} \right) u_{gt} -C^{w}_{gl}w_{g,t+1} && \nonumber\\
& \quad - \left[ C^{w}_{gl}-C^{v}_{gl} \right]^{+} v_{gt} 
&& g\in\mathcal{G}^{1*}, \, t\in \mathcal{T}^0, \, l\in\mathcal{L}_g \label{Staircase9}\\
%
& \delta_{gts} \leq \sum_{i=\underline{T}_{g,s}}^{{\underline{T}_{g,s+1}-1} } w_{g,t-i}
&& g\in\mathcal{G},t\in\mathcal{T}, s\in [1,|\mathcal{S}_{g}|)
%s\in\mathcal{S}_{g} \setminus \left\lbrace S_{g} \right\rbrace
 \label{Startcost1} \\
%
& v_{gt}=\sum_{s=1}^{ |\mathcal{S}_{g}| }\delta_{gts}
&& g\in\mathcal{G},t\in\mathcal{T} \label{Startcost2}\\
%
& c_{gt}^{\text{SU}}=\sum_{s=1}^{|\mathcal{S}_{g}|} C_{g}^{\text{S}} \delta_{gts}
&& g\in\mathcal{G},t\in\mathcal{T} \label{Startcost3}\\
%
& \delta_{gts}=0
&&  g\in\mathcal{G}, s\in [1,|\mathcal{S}_{g}|), \nonumber \\ 
& && t\in\left(\underline{T}_{g,s+1}-DT^{0}_{g}, \underline{T}_{g,s+1} \right) \label{Startcost4}\\
%
& c_{gt}^{\text{SD}}= C_{g}^{\text{SD}} w_{gt} 
&& g\in\mathcal{G},t\in\mathcal{T} \label{Shutdowncost}\\
%
& \sum_{g\in \mathcal{G}} p_{gt}=\textit{De}_{t}
&&  t\in \mathcal{T} \label{System1}\\
%
& \sum_{g\in \mathcal{G}} \overline{p}_{gt} \geq \textit{De}_{t}+R_{t}
&& t\in \mathcal{T} \label{System2}\\
%
& \sum_{g\in \mathcal{G}} r_{gt} \geq R_{t}
&&  t\in \mathcal{T} \label{System3}\\
%
& p_{gt}=p'_{gt}+\underline{P}_{g} u_{gt}
&& g\in\mathcal{G}, t\in\mathcal{T} \label{relation1}\\
%
& \overline{p}_{gt}=\overline{p}'_{gt}+\underline{P}_{g} u_{gt}
&& g\in\mathcal{G},t\in\mathcal{T} \label{relation2}\\
%
& \overline{p}'_{gt}=p'_{gt}+r_{gt}
&& g\in\mathcal{G}, t\in\mathcal{T} \label{relation3}\\
%
& \overline{p}_{gt}=p_{gt}+r_{gt} 
&& g\in\mathcal{G}, t\in\mathcal{T} \label{relation4}\\
%
& p_{gt}\leq \overline{p}_{gt} 
&& g\in\mathcal{G}, t\in\mathcal{T} \label{relation5}\\
%
& p'_{gt}\leq \overline{p}'_{gt} 
&& g\in\mathcal{G}, t\in\mathcal{T} \label{relation6}\\
%
& u_{gt},v_{gt},w_{gt},\delta^{s}_{g}\in \{0, 1\}
&& g\in\mathcal{G}, t\in\mathcal{T}  \label{eq:uvw}\\
%
& p_{gt},p_{gt}',\bar{p}_{gt}',c_{gt}^{\text{p}},c_{gt}^{\text{SU}},r_{gt} \geq 0 
&& g\in\mathcal{G}, t\in\mathcal{T}  \label{eq:pp}
\end{alignat}


The objective function (\ref{eq:obj}) seeks to minimize the total cost, which is composed of the energy production cost $c_{gt}^{\text{p}}$, the fixed cost of operating at a minimum production level $C_{g}^{\text{R}}$. The variable startup cost $c_{gt}^{\text{SU}}$ and  the shutdown cost $c_{gt}^{\text{SD}}$ for each generator $g\in\mathcal{G}$ during a specific time period $t\in\mathcal{T}$. 

%RZ The correspondence of the constraints in the proposed model is listed below, with the names of the original authors.

%RZ \hl{RZ: Antes de modificar esto discutor algo con Uriel}


The minimum uptime and minimum downtime constraints 
are represented by Constraints (\ref{MUT})--(\ref{MDT2}). The logic of the generators' start-up, shut-down, and operation is modeled by Constraints (\ref{garver}). The generation limit requirements are imposed by Constraints (\ref{genlim1})--(\ref{genlim3}).
Constraint (\ref{genlim3}) follows the ramp-up trajectory upper bound of \citet{Pan2016}; the summation terms are
restricted to $t-i\geq 1$, and the shutdown trajectory is represented by the single term in $w_{g,t+1}$ rather than by a
full summation over $T^{\text{RD}}_g$. This is the form used in all the experiments reported in Section \ref{section:Test2};
because (\ref{genlim3}) is a valid inequality, the integer optimum is unaffected by this choice, which only influences the
strength of the linear relaxation. The start and shut-down ramp limits are modeled by Constraints (\ref{startramp})--(\ref{shutdownramp2}). The ramp-up and ramp-down limits are represented by Constraints (\ref{rampup})--(\ref{rampdown}).
%
The staircase production cost are represented by 
Constraints (\ref{Staircase1})--(\ref{Staircase9}). Note that Constraints (\ref{Staircase5}), (\ref{Staircase6}), (\ref{Staircase7}), (\ref{Staircase8}), and (\ref{Staircase9}) are not necessary in the formulation, but they serve to tighten the variables' staircase production.
%
The shut-down cost constraints are Given by Constraints(\ref{Shutdowncost}).

Constraints (\ref{Startcost1})--(\ref{Startcost4}) indicate the cost of the segment of the variable of start-up function cost $c_{gt}^{S_g}$ of the generator $g$, where $C_{g,s}^{\text{S}}$ is the start-up cost in the category $s$ of generator $g$ in \$/MWh.
%
The demand and reserve requirement are modeled by Constraints (\ref{System1}) and (\ref{System3}). Constraints (\ref{relation1})--(\ref{relation6}) establish linear relationships between the power variables 
$p_{gt}, p'_{gt}, \overline{p}'_{gt},\overline{p}_{gt}$. These constraints are used in the best UCP formulations \citep{Knueven2020b}, and help constrain the convex hull of the problem.
%
Lastly, \eqref{eq:uvw} and \eqref{eq:pp}  define the nature of the decision variables.

Initial conditions: Constraints (\ref{garver}),(\ref{rampup}),(\ref{rampdown}) change when the time is $t=0$. The parameters used in this case to substitute the parameter with $t-1$ index are the initial conditions of the generators such as $u_{g,0}, p_{g,0}, p'_{g,0}, U_{g,0}, D_{g,0}$. Also, Constraints (\ref{Startcost4}) assure the start-up variable $\delta_{gt}$ to be zero during the initial periods considering its minimum shut-down time if the generator has been offline \citep{Morales-Espania2013}.







\section{Proposed Matheuristics}
\label{sec:method}

The proposed solution method for the thermal UCP 
consists of two phases: a construction phase, where
an initial feasible solution is built, and an
improvement phase, where five procedures,
including KS and four tailored-made versions of LB,
are applied to the constructed solution.
In the remainder of the section, all these
components are further explained.



\subsection{Construction Phase}
\label{construction}

To the best of our knowledge, the best construction
method for UCP is due to \citet{Harjunkoski2021}.
Their heuristic, referred as HGPS, is of the type relax and fit (R\&F).
The solution strategy of R\&F methods is to reduce the size of the MILP to a smaller, easier-to-solve model called subproblem, fixing those variables likely to keep their values in the optimal solution and letting the solver decide among the others. Finally, the subproblem is input into the solver, expecting that the solver returns a high-quality solution.
Our idea is to develop am enhanced R\&F heuristic that
better exploits the structure of our problem.


%I. Harjunkoski and M. Giuntoli and J. Poland and S. Schmitt
The HGPS heuristic
begins by solving a linear relaxation of the UCP, obtaining a solution $\tilde{x}$. In this method, the variables $u_{g,t}$ to be fixed to one in the subproblem are those that satisfy the condition $\tilde{u}_{g,t}\cdot\tilde{p}_{g,t} \geq \underline{P}$. We will call this condition: Harjunkoski's rule. Let $\tilde{u}_{g,t}$ and $\tilde{p}_{g,t}$ be the values taken by the variables in the linear relaxation solution. Once the variables are fixed, the subproblem is solved, deciding over the rest of the non-fixed variables. 

Our heuristic, named HARDUC,
starts by solving a linear relaxation of the UCP. Then, to form the subproblem, the variables $u_{g,t}$ that satisfy the criterion $\tilde{u}_{g,t}\cdot\tilde{p}_{g,t}=0$ are fixed to zero, and the decision over the other variables is left to the solver.
%\The value of the variables $u_{g,t}$ obtained from the linear relaxation of the problem is represented by $\tilde{u}_{g,t}$.


%\hl{Discuss this paragraph}
Figure \ref{fig:Constructives} illustrate the difference between these two heuristics, indicating the space of ${u}_{g,t}$ variables. On the left-hand side, we see the HGPS method, where the shaded region indicates variables set to one, and the light region shows variables yet to be determined by the solving the subproblem. On the right-hand side, we have the HARDUC method, where the light region comprises non-zero variables according to Harjunkoski's rule but with a value greater than $\underline{P}$. Finally, the variables that meet the criterion ${u}_{g,t}: \{ \tilde{u}_{g,t}\cdot \tilde{p}_{g,t} < \underline{P} \} \cap \{\tilde{u}_{g,t}\cdot\tilde{p}_{g,t}\neq 0\}$ will be used for the construction of a restricted candidate list (RCL). We expect these variables are more likely to be in the optimal solution than those with zero values. This RCL is essential for other improving methods based on local branching described in this work.

\begin{figure}[htb]
\centering
\includegraphics[width=14cm]{Figure/Contructives.pdf}
\scriptsize \caption{Rules for R\&F construction methods. The shaded area represents the subset of variables to be fixed in the subproblem.}
\label{fig:Constructives}
\end{figure}


After obtaining a feasible solution, the search for an improved solution starts as a second phase of the solution algorithm.  We develop a family of five metaheuristics, including one based on KS and four different versions of LB. All these methods are described in the following sections.

\subsection{Improvement Phase}
\label{metaheuristics}

Once an initial feasible solution has been obtained, an improvement phase is applied in order to improve the solution. For this purpose, two main solution frameworks are proposed and explained below.

\subsubsection*{Kernel Search}


Kernel search is a matheuristic initially developed by \citet{Angelelli2010} to solve a multidimensional knapsack problem with outstanding results. 
We have not seen any KS implementation for UCPs, to the best of our knowledge.
%RZ We envision KS as a highly competitive method for solving the UCP, so it was chosen as another matheuristics besides the four LB versions.
%
KS is divided into two phases: initialization and expansion. During the initialization phase, the binary variables in problem $P$ that are likely to take a value of one in the optimal solution are determined. This subset of variables is called the kernel (equivalent to the binary support in LB). The selection of the variables that form the kernel is usually those that take the value one in the linear relaxation.
%
Then, the remaining variables not part of the kernel are sorted according to some economic criterion (often, the reduced costs of the LP relaxation) and divided into small groups called buckets. 

The expansion phase consists of solving the multiple MILP subproblems formed by each bucket concatenated one by one with the kernel, one step at a time. This phase is called expansion because, in each iteration, bucket variables that take the value of one are picked and added to the kernel. The constraint  $\sum_{j\in B_{i}}x_{j} \geq 1$ must be added to the subproblem to enforce that at least one variable of the bucket takes the value one. Because of this constraint, the kernel tends to increase in size with each iteration. For each iteration, the remaining bucket variables that are not part of the subproblem must be fixed to zero. The conventional KS terminates when each bucket has been resolved with the kernel or the allowed time is met \citep{Angelelli2010}. % has timed out 

%\subsubsection*{Implementation of KS for the UCP}
\vspace{1em}\noindent\textit{Implementation of KS for the UCP}\\

We present a version of KS that can be seen in Algorithm \ref{alg:KS} called \texttt{kernelsearch()}. 
The inputs of algorithm \texttt{kernelsearch()} are total time ($\textit{t\_total}$),  a feasible solution ($\bar{x}$) of an instance of UCP to solve ($P$). The outputs are $x^{*}$ and $z^{*}$, respectively, the incumbent solution and its cost.
The function \texttt{kernelsearch()} starts the initialization phase by building the kernel from variables with $u_{g,t}=1$ in the feasible solution $\bar{x}$; this feasible solution is obtained from the constructive method HARDUC. 
It is noted that only the dominant variable $u_{g,t}$ is used for building both the kernel and the buckets. 
Note that our version of KS does not require kernel building. KS is only used as an improvement method, as a constructive method that determines the kernel.

\begin{algorithm}[htb]
\small
\caption{\texttt{kernelsearch()}}
\label{alg:KS}
\begin{algorithmic}
%\fontsize{10}{8}\selectfont 
\Require \\
$\bar{x}$=a feasible solution of $P$, \\
$P$=a formulation of the instance to solve,\\
$\textit{t\_total}$
\Ensure \\
A feasible solution $x^{*}$ of value $z^{*}$ \\

\While{$(elapsed\_time < t\_total)$}
    %\State \hfill
    \State $\textit{cutoff} \leftarrow z^{*}$ \Comment{\textbf{Initialization phase}}
    \State $K \leftarrow$ kernel formed from the variables that fulfill $u_{g,t}:u_{g,t}=1\in\bar{x}$ 
    \State $ P^{ K}\leftarrow$ \texttt{fixing}(1,$K$,$P$) fixing to 1 all the variables into kernel
    \State $\lbrace \hat{x}, \hat{z} \rbrace \leftarrow \texttt{solveLR}( P^{\text{fix(\textit{K})}})$ 
    \State $U \leftarrow$ select the variables that fulfill ${u}_{g,t}\notin K$
    \State $\textit{nbucks} \leftarrow 1 + 3.322 \ln(|\textit{U}|)$
    \State $U^{\text{ord}} \leftarrow$ sort $U$ in ascending order of reduced cost (most negative first) 
    \State $\{ B_i\}_{i=1}^{\textit{nbucks}} \leftarrow$ \texttt{buildbuckets}($U^{\text{ord}}$,$\textit{nbucks}$)  
    \State $tl \leftarrow (\textit{t\_total}-\textit{elapsed\_time}) / n$ 
%RZ    \State \hfill\Comment{\textbf{Expansion phase}}
    \For{($i=1,\textit{nbucks}$)} \Comment{\textbf{Expansion phase}}
            \State $ P^{ K\cup B_i}\leftarrow$ \textbf{fixing}(0,$\overline{B}_i$,$P$) \Comment{fixing to 0 all variables of the other buckets}
            \State $ P^{ K\cup B_i}\leftarrow$ $ P^{ K\cup B_i} \cup \lbrace \sum_{j\in B_{i}}x_{j} \geq 1\rbrace $ \Comment{A variable from the bucket is forced into entering the kernel}
            \State $\lbrace stat, \widetilde{x}, \widetilde{z} \rbrace \leftarrow \texttt{solve}( P^{ K\cup B_i}, \textit{cutoff},tl)$
            \If{(stat = \text{feasible})}
                    \If{$(\widetilde{z} < z^{*} )$} 
                        \State $z^{*} \leftarrow \widetilde{z}$ \Comment{Update best known objective}
                        \State $\textit{cutoff} \leftarrow \widetilde{z}$ \Comment{Tighten cutoff for next iteration}
                        \State $x^* \leftarrow \widetilde{x}$ \Comment{Update best solution}
%                 \State $  \textit{cutoff} \leftarrow z^{*} \leftarrow \widetilde{z}$; $x^* \leftarrow \widetilde{x}$ \hl{ALGO AQUI NO ENCAJA... triple arrow?? UILL - LO CAMBIÉ ARRIBA EN TRES LINEAS EN AZUL Y LO COMENTE PARA ACLARAR CADA LINEA}
                \EndIf
                \State $K \leftarrow$ update the kernel with the variables in $B_i$ such that $u_{g,t} = 1$ in $\widetilde{x}$.
%                \State K $\leftarrow$ update the kernel from the buckets variables that fulfill $\{u_{g,t}:u_{g,t}=1\in\bar{x}\}$ \hl{RZ esta expesion no es correcta UILL doc... la cambie por la de arriba }
            \EndIf
    \EndFor
\EndWhile\\
\Return{$x^{*}$}
\end{algorithmic}
\end{algorithm}          


The construction of the buckets begins by fixing to $1$ the variables constituting the kernel $u_{g,t} \in K$ and solving a linear relaxation of the problem $P^\text{fix(\textit{K})}$.
Then, the set $U$ is constructed with the variables $u_{g,t}$ that are not inside the kernel $K$.
The number of buckets is calculated using Sturges's rule \citep{Scott2009}: $\textit{nbucks}=1+3.322 \ln(|U|)$.

Then, the $u_{g,t}\in U$ variables are sorted in ascending order of their reduced costs obtained from the linear relaxation problem $P^\text{fix(\textit{K})}$, so the most negative reduced costs come first (since this is a minimization problem). The newly ordered set is denoted as $U^{\text{ord}}$.

The set of buckets of size $\textit{nbucks}$, into which the variables of set $U^{\text{ord}}$ are divided, is denoted as $B_i$, where $i$ is an index ranging from 1 to $\textit{nbucks}$.

%https://www.mdpi.com/2076-3417/12/22/11421
\begin{algorithm}[htb]
\small%\fontsize{10}{8}\selectfont 
\caption{\texttt{buildbuckets()}}
\label{alg:buildbuckets}
\begin{algorithmic}
\Require \\
$U^{\text{ord}}$=list of $u_{g,t}$ outside the kernel, most negative reduced cost first, \\
$\textit{nbucks}$=number of buckets
\Ensure \\A set of buckets $\{ B_i\}_{i=1}^{\textit{nbucks}}$\\
\State $\{ B_i\}   \leftarrow \lbrace\phi\rbrace$
\State $\textit{start}      \leftarrow 0 ; \textit{end}  \leftarrow 0$
\State $n                \leftarrow |U^{\text{ord}}|$
\State $k                \leftarrow \lfloor n / \textit{nbucks}\rfloor$
\State $\textit{remainder} \leftarrow n \mod{(\textit{nbucks})}$ 
    \For{ $i=1,\textit{nbucks}$} 
        \State $\textit{end} \leftarrow \textit{start} + k$
        \If{$\textit{remainder} > 0$} 
            \State $ \textit{end}       \leftarrow \textit{end} + 1$
            \State $ \textit{remainder} \leftarrow \textit{remainder} - 1$
        \EndIf            
        \State $ B_i\leftarrow  B_i\cup \lbrace U^{\text{ord}}[\textit{start}:\textit{end}]\rbrace$
        \State $\textit{start} \leftarrow  \textit{end} $
\EndFor\\
\Return{$\{ B_i\}_{i=1}^{\textit{nbucks}}$}
\end{algorithmic}
\end{algorithm}   

Function \texttt{buildbuckets}($U^{\text{ord}}$,$\textit{nbucks}$), depicted in Algorithm \ref{alg:buildbuckets}, is utilized to partition the set of variables $u_{g,t}\in U^{\text{ord}}$ into approximately equal-sized subsets known as buckets $B_i$. First, the size of each bucket is calculated by dividing the total number of variables in $U^{\text{ord}}$ by the desired number of buckets. Also, the remaining variables that cannot be equally distributed are identified. Next, an empty list is initialized to hold the buckets $B_i$, while two variables, $\textit{start}$ and $\textit{end}$, are created to monitor each bucket's range of variables. Then, the function iterates over the number of buckets, sets the range of objects for each subset, and assigns any remainder variables to the first bucket. Next, the range of variables for each bucket is appended to the list $ B_i\leftarrow  B_i\cup \lbrace U^{\text{ord}}[\textit{start}:\textit{end}]\rbrace$. Lastly, \texttt{buildbuckets} returns all the buckets $\{ B_i\}_{i=1}^{\textit{nbucks}}$ as the output.

During the expansion phase of KS, problem $P^{K \cup B_i}$ is solved by concatenating the kernel firstly with the buckets that contain variables having the most negatively reduced costs.
Next, the other buckets' variables $u_{g,t}$ are fixed to zero in the problem, and the constraint $\sum_{j\in B_{i}}x_{j} \geq 1$ is added to ensure that at least one variable in the bucket $B_i$ has a value of one.
The function \texttt{solve}($P$, $\textit{cutoff}$, $\textit{tl}$) is used to send the problem $P^{K\cup B_i}$ to the solver, with $P$ representing the problem, $\textit{cutoff}$ representing the upper bound value, and $tl$ representing the maximum solution time. At each iteration, the kernel is updated, adding the bucket variables that result in a value of one in the solution.

After solving all the buckets $B_i$ with the kernel, if the maximum time has not elapsed, the incumbent solution $x^*$ is set as the new initial solution $\bar{x}$, and the process is restarted until the time is exhausted.



\subsubsection*{Local Branching}

This matheuristic was proposed by \citet{Fischetti2003} and implemented the concept of local search in a MILP problem using a branching strategy. Those feasible solutions within a distance radius of the parameter $k$ define the neighborhood $N(\overline{x},k)$. The distance from solution $\overline{x}$ to other solutions is calculated using the hamming distance $\Delta(x,\overline{x}$), which counts the number of changes from 0 to 1 and from 1 to 0 between the variables that conform the binary support ($\textit{BS}$) of the solution $\overline{x}$ to other variables outside of binary support  ($\overline{BS}$). The binary support of a solution $\overline{x}$ consists of all binary variables that take the value of one in the solution. Exploration in a neighborhood is performed by the solver adding a non-valid inequality called Local Branching Constraint (LBC) (\ref{localbranchingcut}) to problem $P$.

Unlike Fischetti and Lodi, which considers all binary variables of a problem to form a $\textit{BS}$, in this work, the $\textit{BS}$ of a solution $\bar{x}$ of UCP is defined as the subset of the commitment variables {$\textit{BS}$}=$\{u_{g,t}:u_{g,t}=1\}$.
%we define the $\textit{BS}$ of a solution $\bar{x}$ of UCP as the subset of the commitment variables {$\textit{BS}$}=$u_{g,t}:u_{g,t}=1$. 
Therefore, the LBC limits the number of moves between the $\textit{BS}$ and $\overline{\textit{BS}}$ to the number $k$ as defined by the equation:

\begin{align}
&  \Delta(x,\bar{x}) = \sum\limits_{j\in\textit{BS}} (1-u_{j}) + \sum\limits_{j\in{\overline{\textit{BS}}}} u_{j} \leq k \label{localbranchingcut}
\end{align}


In our UCP, the binary variables $v_{g,t}$, $w_{g,t}$, and $\delta_{g,t}$ contained in the formulation are not considered to build the $\textit{BS}$; therefore, they are not used in the local search definition we present. The decision to use only the binary commitment variables $u_{g,t}$ is because we have observed that these variables $u_{g,t}$ are found in most of the model constraints, and a change in the value of these variables can activate or deactivate all the operating constraints related to a generator. Using a dominant variable in the local search has advantages, such as a reduction in the solver search time due to reducing the number of variables in the local search. It also compacts the mathematical model by removing constraints related to a generator. This idea of using a single dominant binary variable was initially reported by \citet{Darwiche2019} and is used in this particular version of LB to solve a graph edit distance problem.



Finally, the constraints representing the movements in LB are enlisted as follows.
\begingroup                 %Group and Split the formulation
\allowdisplaybreaks
\begin{alignat}{2}
&\text{left-branch}: &&  \Delta(x,\bar{x})\leq k \label{eq:leftbranch}\\
& \text{right-branch}: \quad && \Delta(x,\bar{x})\geq k+1 \label{eq:rightbranches}\\
& \text{tabu}:&&  \Delta(x,\bar{x})\geq 1 \label{eq:tabu}\\
& \text{soft-fixing}: &&  \sum\limits_{j\in\textit{BS}}\bar{x}_j x_j \geq 0.9 \sum\limits_{j\in\textit{BS}}\bar{x}_j\label{eq:softfixing}
\end{alignat}
%The procedure described is the standard LB method. 

The constraints in the LB enlisted, as proposed by Fischetti and Lodi in 2006, are crucial for efficiently narrowing the feasible space in solving MILPs. These constraints, which include left-branching, right-branching, tabu constraints, and a soft-fixing strategy, play a key role in the algorithm's effectiveness. This approach, outlined in Fischetti and Lodi's original work, significantly contributes to solution improvement by adjusting the parameter $k$ and implementing diversification mechanisms. Overall, these constraints are essential for the LB method to target near-optimal solutions in MILP problems.

%RZ \subsubsection{Our Implementauon of Local Branching}
\vspace{1em}\noindent\textit{Our implementation of LB for the UCP}\\

This section outlines four versions of LB that differ from the original idea \citep{Fischetti2003}. The main differences are described as follows.
\begin{itemize}
    \item LB1: First, this version narrows the local search between the $\textit{BS}$ and an RCL, which is defined from the elements satisfying the condition ${u}_{g,t}: \{ \tilde{u}_{g,t}\cdot{p}_{g,t} < \underline{P} \} \cap \{\tilde{u}_{g,t}\cdot\tilde{p}_{g,t}\neq 0\}$, a representation of this subset of variables can be shown in the Figure \ref{fig:Constructives}; note that the RCL contains the variables that are not zero but are below the minimum power value of the generator $\overline{P}_{g}$, these variables are discarded by \citet{Harjunkoski2021} in his fixing criterion, but we have included them to construct our RCL.     
    Second, the soft-fixing method further narrows the search, forcing at least 90\% of the original $\textit{BS}$ variables to remain in the solution by applying the constraint 
    $\sum\limits_{j\in\textit{BS}}\bar{x}_j x_j \geq 0.9 \sum\limits_{j\in\textit{BS}}\bar{x}_j$ proposed by \citet{Fischetti2003}. In addition, soft-fixing relaxes the integrality constraint on the $\textit{BS}$ variables but keeps the variables' bounds at [0,1]. Finally, we take advantage of the fact that the binary variables $v_{g,t}$ and $w_{g,t}$ in constraints (\ref{garver}) force $u_{g,t}$ to take binary values, even if $u_{g,t}$ is defined as continuous. The possibility of relaxing the integrality constraint of $u_{g,t}$ and obtaining the same solution without being forcibly binary was visualized and reported by \citet{Morales-Espania2013} and \citet{Morales-Espania2014} in their works about the models T\&C thermal UCP.
    \item LB2: This version is similar to LB1 except that it does not have soft-fixing; therefore, constraints (\ref{eq:softfixing}) are not applied.  %Also, the variables $u_{g,t}$ belonging to $\textit{BS}$ are constrained to a binary domain, meaning they can only take on the values 0 or 1.
    %This version is similar to LB1 except that it does not have soft-fixing; therefore, constraints (\ref{eq:softfixing}) are not applied. Also, the variables $u_{g,t}$ that form the $\textit{BS}$ keep in a binary domain.    
    \item LB3: This version is the closest to the original one proposed by \citet{Fischetti2003}. This version does not do a local search with any RCL and is not narrowed down by soft-fixing. The local search is defined only between the $\textit{BS}$ and the variables that do not form the binary support $\overline{\textit{BS}}$.
    \item LB4: This version is similar to LB1 except that the RCL is formed by the variables with negative values in the reduced costs. The reduced costs are calculated by fixing to 1 of the variables that form the $\textit{BS}$ of the initial solution $\bar{x}$ and solving the linear relaxation of the problem.    
\end{itemize}

A key feature of our versions of LB is that they use only the dominant $u_{g,t}$ variables  to define the $\textit{BS}$, unlike the generic LB of \citet{Fischetti2003}, where all binary variables participate in the local search. %The adaptations to the original LB method are anticipated to expedite the search process, aiming to achieve a solution of increased accuracy.
\section{Computational Experience}
\label{sec:Experimentalwork}

In this section, we present the evaluation
of the proposed algorithms.
All methods were coded in Python version 3.10.0 and the algebraic modeling language Pyomo version 6.6.1 using IBM(R) ILOG(R) CPLEX(R) version 22.1.0.0 as the optimization solver. Tests were carried out in a 64-bit platform with 64 GB of RAM and 2.50 GHz Intel(R) i7(R) 11700 CPU, 65W, on a Linux Ubuntu version 20.0 operating system. 

For all tests the relative optimality gap (ROG) was computed as follows:
%
\begin{equation}
\label{ec:gap}
\text{ROG} = \frac{|LowBound-z\_best|}{((1 \times 10^{-10})+|z\_best|)} \cdot 100 \% \nonumber
\end{equation}
%
where $LowBound$ corresponds to the largest lower (dual) bound obtained by the branch-and-bound solver and $z\_best$, to the best feasible integer solution found by the corresponding method.
%def igap(LB,UB):
%    return abs( LB - UB ) / ( 1e-10 + abs(UB) )    

\subsection{Description of test instances}  
A total of 83 instances, divided into three main subsets (see Table~\ref{tab:instancesKazarlis}), were used to evaluate the proposed methods. The first subset, x7day\_small, consists of 20 instances (10 from \citet{Morales-Espania2013} and 10 derived from an eight-generator dataset proposed by \citet{Ostrowski2012}); while the second (x7day\_medium) and third (x7day\_large)  subsets include 33 and 30 instances, respectively, constructed by gradually increasing the number of generators and randomly combining them to enhance complexity, based on data from \citet{Ostrowski2012}. The planning horizon for all the sets is composed of 168 periods, equivalent to seven days. 
\textcolor{blue}{The configuration of each instance can be consulted in Appendix \ref{app:appD}}

\begin{table}[htb]
\centering
\caption{Description of test instances.}\label{tab:instancesKazarlis}
\begin{tabular}{lrrr}\toprule
Group &Instances &Generators &Instances  \\\cmidrule{1-4}
x7day\_small &20 &28-81 &061-070, 101-110    \\
x7day\_medium &33 &85-156 &071-073, 111-140  \\
x7day\_large &30 &165-405 &074-080, 141-163 \\
\bottomrule
\end{tabular}
\end{table}


The demand profile for each instance, shown in Table \ref{tab:Load}, is obtained by multiplying the demand profile and the sum of the maximum capacity of all generators. A reduction factor of 80\% on a weekday is applied to calculate the weekend demand. The spinning reserve requirement is 5\% of the hourly power demand, except for the 19 instances derived from \citet{Morales-Espania2013} (061--080, excluding 062), which use 10\%.

\begin{table}[H]
\centering
\caption{Load demand profile (\% of total capacity).}\label{tab:Load}
\scriptsize
\begin{tabular}{lrrrrrrrrrrrrr}\toprule
hour &1 &2 &3 &4 &5 &6 &7 &8 &9 &10 &11 &12 \\
demand &71\% &65\% &62\% &60\% &58\% &58\% &60\% &64\% &73\% &80\% &82\% &83\% \\\cmidrule{1-13}
hour &13 &14 &15 &16 &17 &18 &19 &20 &21 &22 &23 &24 \\
demand &82\% &80\% &79\% &79\% &83\% &91\% &90\% &88\% &85\% &84\% &79\% &74\% \\
\bottomrule
\end{tabular}
\end{table}



%%%%%%%%%%%%%%%%%%%%%%%%%%%%%%%%%%%%5
%%%%%%%%%%%%%%%%%%%%%%%%%%%%%%%%%%%%5
%%%%%%%%%%%%%%%%%%%%%%%%%%%%%%%%%%%%%
\subsection{Comparison of constructive methods}
\label{section:Test1}

The goal of this first experiment is to evaluate
the proposed construction heuristic, HARDUC.
To this end, we compare HARDUC, HGPS \citep{Harjunkoski2021}, and the best solution found
by CPLEX (CBS, for CPLEX Best Solution) on all the instances.
%
 CPLEX algorithmic parameters such as \textit{emphasis\_mip}, \textit{mip\_strategy\_file}, and \textit{mip\_tolerances\_mipgap}, were configured as detailed in Table \ref{tab:parameters_solverMILP1}. The maximum computing time was set to 1200 s.  The table shows the average ROG (AROG) and its standard deviation (ROGSD), as well as the number of feasible solutions (NFS) that the each method successfully found.
%Less than two hours is typically the expected time to resolve it \citep{Santos2020}

\begin{table}[H]
\centering
\caption{Parameters of the solver used as constructive CBS.}
\label{tab:parameters_solverMILP1}
\scriptsize
\begin{tabular}{>{\RaggedRight}p{3.5cm}r>{\RaggedRight}p{6cm}r}\toprule
{Parameter} & {Value} & {Description} \\\cmidrule{1-3}
emphasis\_mip &1 &Emphasize feasibility over optimality \\
mip\_strategy\_file &3 &Node file on disk and compressed \\
mip\_tolerances\_mipgap & $1 \times 10^{-6}$  &Relative tolerance between the best integer and the best lower bound \\
\bottomrule
\end{tabular}
\end{table}


Table \ref{tab:descriptive_constructive} shows the comparison of the obtained results. These results indicate that the proposed HARDUC achieved the lowest AROG among all the sets of instances, indicating solutions closer to optimality than CBS and HGPS. It also obtained the lowest ROGSD, ensuring more stable performance. CBS showed the highest AROG and ROGSD in most of the cases. HGPS performed better than CBS but slightly  worse than HARDUC. In terms of the number of feasible solutions, HARDUC successfully found feasible solutions for all the sets of instances. HGPS identified 20/32/29 feasible solutions for the sets of small-, medium-, and large-scale instances, respectively. Finally, CPLEX found 20 feasible solutions in the x7day\_small group (100\%), 26 out of 33 in the x7day\_medium group (78.7\%), and 18 out of 30 in the x7day\_large group (63.3\%).


\begin{table}[H]\centering
\caption{Comparison of constructive methods.}
\label{tab:descriptive_constructive}
\scriptsize
\scalebox{1}{
\begin{tabular}{lrrr|rrr|rrr}\toprule
&\multicolumn{3}{c}{x7day\_small} &\multicolumn{3}{c}{x7day\_medium} &\multicolumn{3}{c}{x7day\_large} \\\cmidrule{2-10}
Method & AROG   & ROGSD  & NFS & AROG   & ROGSD  & NFS & AROG   & ROGSD  & NFS \\\cmidrule{1-10}
CBS    & 0.0075 & 0.0195 & 20  & 0.0005 & 0.0001 & 26  & 0.0071 & 0.0186 & 18 \\
HGPS   & 0.0013 & 0.0003 & 20  & 0.0010 & 0.0002 & 32  & 0.0009 & 0.0004 & 29 \\
HARDUC & 0.0007 & 0.0002 & 20  & 0.0005 & 0.0001 & 33  & 0.0003 & 0.0000 & 30 \\
\bottomrule
%\multicolumn{10}{l}{\footnotesize * Values are in USD. \citep{Morales-Espania2013}}
\end{tabular}
}
\end{table}



Therefore, HARDUC was chosen to obtain an initial feasible solution for the proposed matheuristics evaluated in the next section since it obtained the best performance.

 %In summary, as the instance size increases, the number of solutions found by other methods decreases for , while the proposed constructive method consistently discovers an initial solution within the specified time constraints.



%%% Analysis of variance moved to appendix

%%%%%%%%%%%%%%%%%%%%%%%%%%%%%%%%%%%%%%%%%%%%%
%%%%%%%%%%%%%%%%%%%%%%%%%%%%%%%%%%%%%%%%%%%%%
%%%%%%%%%%%%%%%%%%%%%%%%%%%%%%%%%%%%%%%%%%%%%
\subsection{Evaluation of metaheuristics performance}
\label{section:Test2}

In the following tests, the proposed methods (KS and the four versions of LB) are evaluated under two different computation times: 4000 s.\ and 7200 s.\ (including the time for obtaining the initial feasible solution). For all LB versions, each iteration has a time limit of 1200 s., except the first
subproblem of the search and each subproblem solved right after a diversification restart, which are
limited to 100 s.\ in order to probe the neighborhood cheaply before committing the remaining budget.


\textcolor{blue}{
%Before presenting the comparison with the off-the-shelf solver, 
It is important to clarify that the proposed matheuristics do not operate independently from the solver. Instead, they work alongside CPLEX, invoking the solver as an exact component repeatedly throughout the search. In this sense, the matheuristics function in tandem with the solver, guiding and accelerating its convergence rather than replacing it. To assess the extent of this synergy, we also include experiments in which the model is solved solely by CPLEX, allowing us to verify whether the combined methods indeed outperform the solver on its own.
}

\color{blue}
%\subsubsection{Comparison with an off-the-shelf solver} 

To establish this baseline, we evaluate the performance of CPLEX when solving the TUCP model \eqref{eq:obj}–\eqref{eq:pp}. Two benchmark methods were considered. In Solver Method 1 (SM1), the TUCP model is solved directly by CPLEX without any warm start. In Solver Method 2 (SM2), an initial feasible solution generated by our constructive heuristic is provided to the solver as a MIP start. For all methods (except SM1), the same initial solution is used to ensure a fair comparison. The specific CPLEX parameters adopted for these experiments are reported in Table~\ref{tab:parameters_solverMILP2}.
\color{black}

%the results obtained by solving model TUCP ( \eqref{eq:obj}-\eqref{eq:pp}) using CPLEX, are analyzed. For this purpose, two versions were considered: SM1 (Solver method 1), in which the model is completely solved by CPLEX, and SM2 (Solver method 1), where an initial feasible solution is given as a staring solution (MIPStart) to the solver. For all methods (except SM1), the same initial solution is considered. Certain parameters used by CPLEX were adjusted as specified in Table \ref{tab:parameters_solverMILP2}.

%The solver parameters for these improvement methods have been tuned to emphasize the search for feasible solutions over the search for optimal solutions \citep{Fischetti2003}.

\begin{table}[H]
\centering
\caption{CPLEX algorithmic parameters for matheuristic evaluation.}
\label{tab:parameters_solverMILP2}
\scriptsize
\begin{tabular}{c>{\RaggedRight}p{3.5cm}r>{\RaggedRight}p{6cm}r}\toprule
Method & {Parameter} & {Value} & {Description} \\\cmidrule{1-4}
\multirow{4}{*}{{SM1}} &emphasis\_mip &1 &Emphasize feasibility over optimality \\
&mip\_strategy\_file &3 &Node file on disk and compressed \\
&mip\_tolerances\_mipgap & $1 \times 10^{-6}$ &Relative tolerance between the best integer and the best lower bound \\
&preprocessing\_symmetry &$-1$ &CPLEX default (automatic symmetry detection) \\
\bottomrule \bottomrule
\multirow{4}{*}{{Remaining}} &emphasis\_mip &1 &Emphasize feasibility over optimality \\
&mip\_strategy\_file &3 &Node file on disk and compressed \\
&mip\_tolerances\_mipgap & $1 \times 10^{-6}$ &Relative tolerance between the best integer and the best lower bound \\
&preprocessing\_symmetry &0 &Turn off symmetry breaking \\
\bottomrule
\end{tabular}
\end{table}

Any CPLEX parameter not listed in Table \ref{tab:parameters_solverMILP2} was left at its default value.

For the relevant criteria specific to each proposed solution method, Table \ref{tab:methods} illustrates how these were applied.

\begin{table}[htpb]
\centering
\small
\caption{Criteria considered for the KS and the LB methods.}\label{tab:methods}
\begin{tabular}{p{1.5cm}ll}\toprule
\multirow{3}{*}{LB1} 
& LBC\tnote{1}~built using dominant variable $u_{g,t}$. \\
& RCL from Harjukovski's rule \citep{Harjunkoski2021}. \\
& Soft-fixing to 90\% of binary support. \\ 
\cmidrule{2-2}
\multirow{2}{*}{LB2} 
& LBC built using dominant variable $u_{g,t}$. \\
& RCL from Harjukovski's rule \citep{Harjunkoski2021}. \\ 
\cmidrule{2-2}
\multirow{1}{*}{LB3} 
& LBC built using dominant variable $u_{g,t}$. \\  
\cmidrule{2-2}
\multirow{3}{*}{LB4} 
& LBC built using dominant variable $u_{g,t}$. \\
& RCL from reduced costs. \\ 
& Soft-fixing to 90\% of binary support. \\ 
\cmidrule{2-2}
\multirow{4}{*}{KS} 
%& The first solution is given by the HARDUC constructive method. \\
& Kernel and buckets are built only with dominant variables $u_{g,t}$. \\
& Number of buckets is calculated by Sturge's rule \citep{Scott2009}. \\
& Buckets are built using the reduced costs from linear relaxation (fixing kernel). \\
\bottomrule
\end{tabular}
\end{table}
%&The buckets are built using the ordered reduced costs fixing the Kernel and solving linear relaxation. \\

Table \ref{tab:descriptive_fused} shows the AROG, the ROGSD, and the NFS obtained for each method considering the set of test instances. As can be seen, on average, increasing the maximum computation time reduce the average ROG and its corresponding standard deviation in most cases. Considering the significant impact that improvements can have, even if they are perceived to be small, we can visualize that by assigning one more hour of optimization (which is still within the time allowed in the real case), we can further improve the solutions obtained when the optimization time is 4000 s.

%limit. However, some medium- (NFS = 26/33) and large-scale (NFS = 28/30) instances were excluded in the 

Our proposed construction heuristic, HARDUC, found feasible solutions for all instances within the time limit. However, some medium- and large-scale instances were excluded in the computations under the 7200 s. \textcolor{blue}{option because the base solver failed to find an initial feasible integer solution at the root node of the branch-and-bound method, so no upper bound was available, although the linear relaxation (lower bound) could be obtained. We point out in Appendix~\ref{app:appD} the instances without an initial integer solution. This issue (failure to find an initial feasible integer solution) has also been reported in previous studies \citep{Cong2015,Knueven2020b,Morales-Espania2015,Wuijts2024} }.

%WRONG!!! option because CPLEX, the base solver, failed to solve the linear programming relaxation at the root node of the branch-and-bound method so that no lower bound was available. !!!



%solve them in 7000 seconds, and therefore statistics for SM! cannot eb computed in those cases. %As the ROG metric depends on the best solution found, these cases were omitted. TENGO DUDAS!!}



\begin{table}[htpb]
\centering
\caption{Performance comparison among all solution methods considering different maximum time limits}
\label{tab:descriptive_fused}
\scriptsize
\begin{tabular}{llrrr|rrr}
\toprule
\cmidrule{2-8}
\multicolumn{1}{c}{Set} &Method &\multicolumn{3}{c}{4000 s} & \multicolumn{3}{c}{7200 s}\\
\cmidrule{3-8}
& & AROG & ROGSD & NFS & AROG & ROGSD & NFS \\
\midrule
 \multirow{7}{*}{x7day\_small} & LB1 &6.65E-4 &2.20E-4 &20 &6.11E-4 &1.82E-4 &20 \\
&LB2 &  6.56E-4 &2.18E-4 &20 &6.22E-4 &2.00E-4 &20 \\
&LB3 &  6.98E-4 &2.44E-4 &20 &6.67E-4 &2.31E-4 &20 \\
&LB4 &  6.60E-4 &2.30E-4 &20 &6.31E-4 &2.06E-4 &20 \\
&KS  &  6.85E-4 &2.63E-4 &20 &6.78E-4 &2.67E-4 &20 \\
&SM1 &  5.79E-4 &1.59E-4 &20 &5.56E-4 &1.59E-4 &20 \\
&SM2 &  6.28E-4 &1.49E-4 &20 &5.81E-4 &1.55E-4 &20 \\
\midrule
\multirow{7}{*}{x7day\_medium} &LB1  &3.74E-4 &9.08E-5 &26 &3.36E-4 &7.25E-5 &26 \\
&LB2  &3.66E-4 &9.49E-5 &26 &3.39E-4 &7.67E-5 &26 \\
&LB3  &3.82E-4 &9.44E-5 &26 &3.50E-4 &6.46E-5 &26 \\
&LB4  &3.76E-4 &9.19E-5 &26 &3.38E-4 &6.91E-5 &26 \\
&KS   &3.10E-4 &7.48E-5 &26 &2.98E-4 &6.73E-5 &26 \\
&SM1  &4.73E-4 &1.18E-4 &26 &3.13E-4 &8.77E-5 &26 \\
&SM2  &3.74E-4 &9.08E-5 &26 &4.45E-4 &1.07E-4 &26 \\
\midrule
\multirow{7}{*}{x7day\_large} & LB1  &2.71E-4 &6.42E-5 &18 &2.47E-4 &5.75E-5 &18 \\
&LB2 &  2.81E-4 &7.17E-5 &18 &2.42E-4 &5.82E-5 &18 \\
&LB3 &  2.91E-4 &7.36E-5 &18 &2.63E-4 &6.44E-5 &18 \\
&LB4 &  2.90E-4 &6.64E-5 &18 &2.54E-4 &6.30E-5 &18 \\
&KS  &  2.25E-4 &3.95E-5 &18 &2.08E-4 &3.82E-5 &18 \\
&SM1 &  3.07E-4 &8.77E-5 &18 &2.72E-4 &7.33E-5 &18 \\
&SM2 &  3.33E-4 &7.79E-5 &18 &3.19E-4 &8.34E-5 &18 \\
\bottomrule
%\multicolumn{8}{l}{\footnotesize * Values are in USD. \citep{Morales-Espania2013}}
\end{tabular}
\end{table}


Finally, the results are also plotted for all methods and each set of instances. A statistical analysis is conducted only for those instances in which the method successfully found a solution. Figures \ref{fig_violinAll_1h_x7day_small}, \ref{fig_violinAll_1h_x7day_medium}, and \ref{fig_violinAll_1h_x7day_large} present the distributions of the relative optimality gap for each method. Each series represents a method, with the horizontal axis indicating the relative optimality gap. The Y-axis represents the density, which is a continuous estimate of the relative frequency. Therefore, the area under the curve must equal 1. The figures show that as instance difficulty increases, the KS and LB methods yield results with significantly lower variance compared to the solver (SM1, SM2). Additionally, the KS and LB methods achieve higher accuracy, with an average relative optimality gap closer to zero than the solver. To assess the significance of these improvements, we conducted statistical tests by analyzing variance and comparing means. 

Similarly, Figures \ref{fig_violinAll_2h_x7day_small}, \ref{fig_violinAll_2h_x7day_medium}, and \ref{fig_violinAll_2h_x7day_large} provide evidence of the relative optimality gap distributions for each method at 7200 s. It can be noticed that, as instance difficulty increases, the KS outperforms the other methods, showing a lower variance and a superior accuracy of AROG.

An analysis of variance (ANOVA) was conducted for each set of instances in both cases. The results reject the null hypothesis of equal means, confirming that at least one method differs significantly from the others. Similar results are obtained when the maximum time limit allowed is 7200 s (see Appendix \ref{app:appB} for details).

%\hl{Caray, los gaps son tan pequenos que tenemos que decir algo donde la escala impacte.  Por ej.}
%\hl{It is important to note that even though the relative optimality gaps are in the order of 1E-4, the dimension of the objective function value is in the order of TENS of billions of dollars, thus, even a small improvement of 0.01\% or 0.001\% is meaningful in practice.  Es cierto esto? UILL - AGREGUE LOS PARRAFOS  DE ABAJO EN AZUL  PARA ACLARAR ESTO: SE ACEPTAN COMENTARIOS/ MEJORAS }


It is important to realize that even very small relative optimality gaps, on the order of 0.01\% or less, can lead to materially different market outcomes. These differences may not significantly affect the total cost, but they can shift unit commitment decisions enough to alter the marginal prices used for market settlements. Since these prices determine generator payments, slight variations can result in economic mismatches and fairness concerns \citep{Johnson1997}.

This sensitivity has been confirmed in several studies, where near-optimal solutions with nearly identical objective values led to significantly different pricing outcomes. As shown by \citet{Sioshansi2008}, small optimality gaps do not necessarily eliminate payment deviations or price volatility. Therefore, reducing the gap (even marginally) remains relevant for achieving more stable and economically efficient market results. In this work, we aim to contribute to reducing the optimality gap as a means to help mitigate, among other issues, the payment instability faced by generators.

%%%%%%%%%%%%%%%%%%%%%%%%%%%%%%%%%%%%%%% x7day_small 
\begin{figure}[htb]
     \centering
     \begin{subfigure}[b]{0.48\textwidth}
         \centering
         \includegraphics[width=\textwidth]{Figure/fig_violinAll_1h_x7day_small.pdf}
         \caption{Time limit of 4000 seconds}
         \label{fig_violinAll_1h_x7day_small}
     \end{subfigure}
     \hfill
     \begin{subfigure}[b]{0.48\textwidth}
         \centering
         \includegraphics[width=\textwidth]{Figure/fig_violinAll_x7day_small.pdf}
         \caption{Time limit of 7200 seconds}
         \label{fig_violinAll_2h_x7day_small}
     \end{subfigure}
     \caption{Relative optimality gap distribution of improvement methods on set x7day\_small.}
\end{figure}



%%%%%%%%%%%%%%%%%%%%%%%%%%%%%%%%%%%%%%%% x7day_medium 
\begin{figure}[htpb]
     \centering
     \begin{subfigure}[b]{0.485\textwidth}
         \centering
         \includegraphics[width=\textwidth]{Figure/fig_violinAll_1h_x7day_medium.pdf}
         \caption{Time limit of 4000 seconds}
         \label{fig_violinAll_1h_x7day_medium}
     \end{subfigure}
     \hfill
     \begin{subfigure}[b]{0.485\textwidth}
         \centering
         \includegraphics[width=\textwidth]{Figure/fig_violinAll_x7day_medium.pdf}
         \caption{Time limit of 7200 seconds}
         \label{fig_violinAll_2h_x7day_medium}
     \end{subfigure}
     \caption{Relative optimality gap distribution of improvement methods on set x7day\_medium.}
\end{figure}

%%%%%%%%%%%%%%%%%%%%%%%%%%%%%%%%%%%%%%%%%%%%%% x7day_large 
\begin{figure}[htb]
     \centering
     \begin{subfigure}[b]{0.48\textwidth}
         \centering
         \includegraphics[width=\textwidth]{Figure/fig_violinAll_1h_x7day_large.pdf}
         \caption{Time limit of 4000 seconds}
         \label{fig_violinAll_1h_x7day_large}
     \end{subfigure}
     \hfill
     \begin{subfigure}[b]{0.48\textwidth}
         \centering
         \includegraphics[width=\textwidth]{Figure/fig_violinAll_x7day_large.pdf}
         \caption{Time limit of 7200 seconds}
         \label{fig_violinAll_2h_x7day_large}
     \end{subfigure}
     \caption{Relative optimality gap distribution of improvement methods on set  x7day\_large.}
\end{figure}

\clearpage

%\section{Discussion}
\subsection{Evaluation of special cases}
\label{sec:Discussion}

To gain deeper insight, a more detailed analysis of the following cases is carried out as follows:
\begin{description}
\itemsep0.05em
    \item[Case 1:] Performance of methods in two instances where SM1 was unable to find a solution
    \item[Case 2:] Performance of methods on two instances where SM1 found a solution but performed worse than the remaining methods
    \item[Case 3:] Performance of methods on two instances for which SM1 obtained similar results to the remaining methods
    \item[Case 4:] Comparison of the methods in instances where the KS shows stagnation.
    \item[Case 5:] Comparison of methods in which KS failed to improve the initial solution
\end{description}


\subsubsection*{Analysis of Case 1}

In this case two instances,uc\_141 and uc\_146 from the set of large-scale instances, are evaluated in which SM1 was not able to find a feasible solution. As can be seen in Figures \ref{Fig_uc_141_x7day_largeA} and \ref{Fig_uc_146_x7day_largeA}, the proposed solution methods (and their variants) were able to improve the initial solution found by HARDUC, with KS performing the best of all and LB3 showing the worst performance.

\begin{figure}[H]
     \centering
     \begin{subfigure}[b]{0.45\textwidth}         \centering
         \includegraphics[width=\textwidth]{Figure/Fig_uc_141_x7day_large.pdf}
         \caption{Instance 141.}            \label{Fig_uc_141_x7day_largeA}
     \end{subfigure}
     \begin{subfigure}[b]{0.45\textwidth}         \centering
         \includegraphics[width=\textwidth]{Figure/Fig_uc_146_x7day_large.pdf}
         \caption{Instance 146.}            \label{Fig_uc_146_x7day_largeA}
     \end{subfigure}
        \caption{Comparison of methods on instances where the solver (SM1) could not solve within the allowed time limit.} \label{fig:Fig_uc_141146_x7day_large}
\end{figure}

\subsubsection*{Analysis of Case 2}

In the second case, two instances (uc\_143 and uc\_153), depiucted in Figures \ref{Fig_uc_143_x7day_largeA} and \ref{Fig_uc_153_x7day_largeA}, illustrate a scenario where SM1 finds a solution, but its performance and quality are not competitive with the other methods. Here, the initial solution provided by HARDUC plays a crucial role in both the starting point and overall performance of the proposed methods. Although SM1 required a significant amount of time to achieve its first notable improvements, the other methods have taken advantage of the initial solution of HARDUC much earlier, allowing them to use the remaining time more effectively for refinement. Among these methods, KS again demonstrated superior performance and solution quality.

\begin{figure}[H]
     \centering
     \begin{subfigure}[b]{0.45\textwidth}         \centering
         \includegraphics[width=\textwidth]{Figure/Fig_uc_143_x7day_large.pdf}
         \caption{Instance 143.}            \label{Fig_uc_143_x7day_largeA}
     \end{subfigure}
     \begin{subfigure}[b]{0.45\textwidth}         \centering
         \includegraphics[width=\textwidth]{Figure/Fig_uc_153_x7day_large.pdf}
         \caption{Instance 153.}            \label{Fig_uc_153_x7day_largeA}
     \end{subfigure}
        \caption{Performance of methods on two instances where SM1 found a solution but performed worse than the remaining methods} \label{fig:Fig_uc_143153_x7day_large}
\end{figure}

\subsubsection*{Analysis of Case 3}

In some isolated cases, SM1 produced competitive or even superior results compared to metaheuristic methods. As shown in Figures \ref{Fig_uc_074_x7day_largeA} and \ref{Fig_uc_079_x7day_largeA}, the initial solution of SM1 was significantly better than the starting point provided by HARDUC. However, while the proposed methods quickly achieved substantial improvements in most cases, they were unable to surpass the solution of SM1 in these specific instances. Only the KS was able to arrive at a similar solution in the given time in one of the cases, being in general one of the proposed methods that showed a better performance.  In particular, despite the strong initial solution of SM1, it struggled to make significant further improvements for most of the remaining time.

\begin{figure}[H]
     \centering
     \begin{subfigure}[b]{0.45\textwidth}         \centering
         \includegraphics[width=\textwidth]{Figure/Fig_uc_074_x7day_large.pdf}
         \caption{Instance 074.}            \label{Fig_uc_074_x7day_largeA}
     \end{subfigure}
     \begin{subfigure}[b]{0.45\textwidth}         \centering
         \includegraphics[width=\textwidth]{Figure/Fig_uc_079_x7day_large.pdf}
         \caption{Instance 079.}            \label{Fig_uc_079_x7day_largeA}
     \end{subfigure}
        \caption{Performance of methods on two instances for which SM1 obtained similar results to the remaining methods} \label{fig:Fig_uc_074079_x7day_large}
\end{figure}


\subsubsection*{Analysis of Case 4}

In Case 4, two instances show a particular pattern in the KS method. Initially, KS achieved significant improvements, but as iterations progressed, it reached longer periods of stagnation, where further enhancements became minimal. This stagnation prevented the algorithm from making meaningful progress in subsequent iterations. For example, in Figure \ref{Fig_uc_075_x7day_largeA}, the solution improved substantially during the first iteration, but showed little progress over the remaining 5000 seconds. A similar pattern is observed in Figure \ref{Fig_uc_161_x7day_largeA}, where although the solution improved considerably in the final 1000 seconds, it struggled for most of the runtime to escape a local optimum. This behavior may be due to the first buckets containing key variables that provide early improvements, while the remaining buckets, due to their size, failed to capture the ideal combination of variables needed for further optimization. It would be valuable to explore whether evaluating a larger kernel or expanding the number of buckets could improve efficiency and mitigate stagnation.

\begin{figure}[H]
     \centering
     \begin{subfigure}[b]{0.45\textwidth}         \centering
         \includegraphics[width=\textwidth]{Figure/Fig_uc_075_x7day_large.pdf}
         \caption{Instance 075.}            \label{Fig_uc_075_x7day_largeA}
     \end{subfigure}
     \begin{subfigure}[b]{0.45\textwidth}         \centering
         \includegraphics[width=\textwidth]{Figure/Fig_uc_161_x7day_large.pdf}
         \caption{Instance 161.}            \label{Fig_uc_161_x7day_largeA}
     \end{subfigure}
        \caption{Comparison of the methods in instances where the KS shows stagnation} \label{fig:Fig_uc_075161_x7day_large}
\end{figure}


\subsubsection*{Analysis of Case 5}

In this final analysis, there are four cases where KS failed to improve the initial feasible solution provided by HARDUC (see Figures \ref{Fig_uc_155_x7day_largeA}, \ref{Fig_uc_156_x7day_largeA}, \ref{Fig_uc_158_x7day_largeA}, and \ref{Fig_uc_159_x7day_largeA}). In these cases, only the LB method and its variations were able to achieve better solutions, with slightly different performance depending on the stage solved. One possible approach to address this limitation is to analyze alternative kernel and hub sizes to determine whether tuning these parameters could increase the opportunity for further improvements.

\begin{figure}[htb]
     \centering
     \begin{subfigure}[b]{0.45\textwidth}         \centering
         \includegraphics[width=\textwidth]{Figure/Fig_uc_155_x7day_large.pdf}
         \caption{Instance 155.}            \label{Fig_uc_155_x7day_largeA}
     \end{subfigure}
     \begin{subfigure}[b]{0.45\textwidth}         \centering
         \includegraphics[width=\textwidth]{Figure/Fig_uc_156_x7day_large.pdf}
         \caption{Instance 156.}            \label{Fig_uc_156_x7day_largeA}
     \end{subfigure}
     \begin{subfigure}[b]{0.45\textwidth}         \centering
         \includegraphics[width=\textwidth]{Figure/Fig_uc_158_x7day_large.pdf}
         \caption{Instance 158.}            \label{Fig_uc_158_x7day_largeA}
     \end{subfigure}
     \begin{subfigure}[b]{0.45\textwidth}         \centering
         \includegraphics[width=\textwidth]{Figure/Fig_uc_159_x7day_large.pdf}
         \caption{Instance 159.}            \label{Fig_uc_159_x7day_largeA}
     \end{subfigure}
        \caption{Comparison of methods on instances where the KS remained in the initial solution without improvement.} \label{fig:Fig_uc_156158159_x7day_large}
\end{figure}


Briefly speaking, since the solver did not find a feasible solution in some large-scale instances, we provided an initial starting solution using our HARDUC method. However, even with initial solution, the solver obtained less accurate results than the proposed improvement methods LB1-4, and KS.

For small-scale instances, when the maximum time limit was set to 4000 s., the solver performed similarly to the LB1-4 and KS methods, with no significant accuracy. However, when the time was extended to 7200 s., the solver outperformed LB3 and SM2, achieving a smaller relative optimality gap.

For medium-scale instances with a 4000-second time limit, the KS outperformed the remaining methods. In addition, LB1-4 performed better than SM1 and SM2.

For large-scale instances, in general, KS was again the most efficient method for both time limits. In addition, LB1 performed better than SM1, and SM2 outperformed SM1. 

It was observed that all LB methods outperformed SM2. However, when the time limit was set to 7200 s., no significant differences were found between LB1-4 and SM1 methods. It is important to note that the solver did not successfully solve all medium- and large-scale instances. Furthermore, if an initial solution is provided to the solver, the performance is noticeably lower compared when it is performed without it.


In general, among the methods, KS achieved the smallest average optimality gap. On the other hand, LB3, which closely follows the original method proposed by \citet{Fischetti2003}, has shown the lowest performance among the proposed LB methods. 

Although KS generally provided better results in terms of relative optimality gap, it did not improve the initial HARDUC solution in approximately 4\% of the instances. In contrast with the LB methods, which successfully improved the initial solution.

\section{Conclusions}
\label{sec:Conclusions}

Five methods were developed to solve a thermal UCP within a matheuristic approach, four based on local branching and one based on kernel search. In addition, a constructive heuristic (HARDUC) was developed to provide the first solution to the matheuristic methods.

Among the methods based on LB, LB3 is the closest version to the original local branching, unlike LB1, LB2, and LB4, which are variants that implement the soft-fixing concept and a restricted candidate list. Additionally, only the variables classified as dominant are considered in the local search. In this case, the commitment variable, which determines the ``on'' or ``off'' state of the generator, was identified as such.

%%%%%%%%%%%%%%%

The methods were assessed by solving three groups of instances classified into small-, medium-, and large-scale according to the number of generators they contained. A time limit of 4000 seconds and 7200 seconds was set for solving the instances using the LB1, LB2, LB3, LB3, LB4, and KS improvement methods, including the commercial solver. We also tried out feeding the off-the-shelf solver with a feasible solution
found by the constuction heuristic.  
%The results of the methods were compared against the best solution achieved by the solver providing it with an initial feasible solution and without an initial feasible solution. 
%The test instances were created based on existing instances from the literature that are known to be challenging to solve.

%%%%%% 
The solver without an initial solution only solved approximately 79\% of medium-sized instances and 63\% of large-scale instances. In contrast, our methods, LB1, LB2, LB3, LB4, and KS, always had feasible initial solutions provided by our HARDUC method in the constructive phase.

In terms of solution quality, statistical tests showed no significant differences among the LB1, LB2, LB4, LB3, and KS methods and the solver on small-scale instances. However,
for testing medium-scale instances within a maximum allowed running time of 4000 seconds, the KS achieved the best performance with statistical significance. The other methods outperformed the solver with an initial solution and matched the results of the solver without an initial solution. For tests with a maximum allowed running time of 7200 s., the solver performed in the same way as all methods except for LB3, which performed worse than the solver without an initial solution. 
LB3 is the method most similar to the original version of local branching.
%
Regardless of the time limit, all methods outperformed the solver in solution quality for the set of large-scale instances.
Again, the KS performed best.

Therefore, we can confirm that the proposed methods are more efficient than using only the solver; moreover, KS is the best method for medium- and large-scale instances.
According to the results, the proposed adaptations of local branching, including soft-fixing and RCL components, helped find better solutions in medium-scale instances instead of the original version.

Although previous studies have suggested that the KS method effectively solves knapsack problems with promising outcomes, the results of our study demonstrate that the KS method may also be useful in solving the UCP problem.

The KS exhibits a behavior similar to a greedy algorithm, with a quick descent followed by a ``stagnation'' effect. However, local branching methods offer a consistent improvement and have the capacity to avoid local optima, albeit with a slower descent rate than KS. A promising direction for future work would be to hybridize KS with LB to leverage the strengths of both methods.

We have learned that matheuristics are able to ``dive'' within the solution space and help speed up the search for better solutions than if we only used the solver. 
%
We have found that the proposed improvement methods do their job by improving the initial solution. Therefore, we can consider our strategy of constructing and improving solutions using matheuristic methods to be effective.

Finally, in practice, we recommend using these matheuristic methods and the solver simultaneously on different computers. This strategy of diversifying solution methods always allows us to obtain feasible and high-quality solutions.

\noindent\textit{Future research}:
For future research, it would interesting the developing of a hybrid algorithm combining KS and the Local Branching Constraint (LBC) (\ref{localbranchingcut}) to benefit from KS's rapid convergence in initial iterations and LBC's ability to escape local optima when KS stagnates. Another
line of work could be designing a branch-and-cut method that includes local branching constraints, using callback functions to introduce cuts into the solver's solution tree, saving time and enhancing efficiency during problem solving.
\textcolor{blue}{We also consider it worthwhile to assess the proposed approach across various UCP models, including variations such as prohibited zones, non-linear elements, elastic demand, interdependent generators, and network flow limits. Incorporating a full network representation, either a linear DC power flow model or a conic AC power flow, increases computational complexity and further motivates advanced solution methods, while the proposed matheuristics can be integrated into a Security-Constrained Unit Commitment (SCUC) framework with minimal structural modifications.}
%
%
Another line of work, within the context of KS, consists of further studies about strategies for determining values for the kernel and bucket list size. \textcolor{blue}{Another avenue is to construct either the kernel in the KS or the binary support in Local Branching using variables identified through learning-based methods as having a high probability of assignment.} 
%
Naturally, many of the ideas developed here can be extended to handle stochastic versions of the problem.  Many of those solution algorithms rely on decomposition or scenario-based schemes where some subproblems might have similar structures.

%\textcolor{blue}{DH-Agregar: evaluar diferentes estrategias para determinar el tamaño de kernel y buckets debido a que en algunas instancias no fue posible mejorar la solución en diversas iteraciones o no se pudo mejorar la primera solución factible encontrada por el KS. También analizar otras estrategias para determinar la mejor combinación de parámetros a fijar para las versiones de LB.}

\vspace{1em}
\noindent
\textit{Acknowledgments:}
\textcolor{blue}{We are grateful to two anonymous reviewers whose
criticism helped improve the quality of this work.}
Uriel Lezama-Lope acknowledges a scholarship for doctoral studies from the Mexican Council for Humanities, Science, and Technology (CONAHCYT).
The research of Roger R\'{\i}os-Mercado was funded
by CONAHCYT under grant CF-2023-I-880.
The research of Diana Huerta-Mu\~noz has been supported by CONAHCYT through a postdoctoral fellowship.

%\section{Data Availability}
%\label{sec:DataAvailability}
\vspace{1em}
\noindent\textit{Data Availability:} The repository \url{https://github.com/urieliram/tc_uc} contains the data and source code for the methods that have been implemented.





\bibliographystyle{plainnat} 
\bibliography{refs_ucp}

\newpage
\appendix
\setcounter{table}{0}
\setcounter{figure}{0}
\renewcommand{\thetable}{A\arabic{table}}
\renewcommand{\thefigure}{A\arabic{figure}}

%RZ \section*{Appendix}\label{app:appendix}

%RZ\subsection*{Appendix A} \label{app:appA}
\section{Results for Construction Heuristics} \label{app:appA}


%\hl{Las figs no se les ve le leyenda del eje x, es decir, se ven empalmados los valores, (UILL ya quedo arreglado...)}

The distributions of the relative optimality gap results for each method are shown in Figures \ref{fig:violinHARDUC_x7day_small}, \ref{fig:violinHARDUC_x7day_medium}, and \ref{fig:violinHARDUC_x7day_large}. Each series represents the results of a method; the relative optimality gap is shown on the horizontal axis. The number of feasible instances found from each group's total instances is reported in parentheses. 
\textcolor{black}{The figures illustrate that, as the instance difficulty increases, the relative optimality gap of the first solution obtained with the HARDUC method exhibits a distribution with significantly lower variance than other methods. Moreover, it demonstrates superior accuracy, as evidenced by the average relative optimality gap approaching zero. Statistical tests further support these findings.}

%%%%%%%%%%%%%%%%%%%%%%%%%%%%%%%%%%%%%%%%%%%%% HGPS VS HARDUC %%%%%%%%%%%%%%%%%%%%%%%%%%%%%%%%%%%%%%%%%%%
\begin{figure}[htb]
\centering
\includegraphics[width=.7\textwidth]{Figure/fig_violinHARDUC_x7day_small.pdf}
\scriptsize \caption{Relative optimality gap distributions of constructive methods in the instances group x7day\_small.}\label{fig:violinHARDUC_x7day_small}
\end{figure}

\begin{figure}[htb]
\centering
\includegraphics[width=.7\textwidth]{Figure/fig_violinHARDUC_x7day_medium.pdf}
\scriptsize \caption{Relative optimality gap distributions of constructive methods in the instances group x7day\_medium.}\label{fig:violinHARDUC_x7day_medium}
\end{figure}

\begin{figure}[htb]
\centering
\includegraphics[width=.7\textwidth]{Figure/fig_violinHARDUC_x7day_large.pdf}
\scriptsize \caption{Relative optimality gap distributions of constructive methods in the instances group x7day\_large.}\label{fig:violinHARDUC_x7day_large}
\end{figure}

The following is an analysis of variance, and mean difference applied to the results of the methods regarding the relative optimality gap. 

On one hand, the analysis of variance checks whether the average relative optimality gap of the results of all methods is statistically different from each other. 
On the other hand, a series of mean difference tests applied to each pair of results verify that the average relative optimality gap of one method is significantly higher than the mean obtained by another method. 
All statistical studies were done only in the instances where the solver found a feasible solution.


The results of the analysis of variance applied to the results of the constructive methods are shown in Table \ref{tabla:ANOVA_constructive}. In the three groups of instances, the null hypothesis $\text{H}_o$ of equality of means was rejected. The alternative hypothesis that at least one of the methods has a mean difference from the others was not rejected.  
The results of the mean difference analysis between the HARDUC and HGPS methods can be found in Table \ref{tabla:Differences_constructive}. The HARDUC method has a significantly lower average in the relative optimality gap than HGPS.

Because the HARDUC method obtained a feasible initial solution in all test instances and a significantly smaller relative optimality gap than those obtained by HGPS and CBS, we consider the HARDUC constructive method to be the most appropriate to provide the first initial solution to the improvement methods (LB1, LB2, LB3, LB4, and KS) in the following tests. 

%%%%%%%%%%%%%%%%%%%%%%%%%%%%%%%%%%%%%%%%%%%%%%%%%%%%%%
\begin{table}[H]\centering
\caption{Analysis of variance summary of the constructive methods HARDUC, HGPS, and CBS.}\label{tabla:ANOVA_constructive}
\scriptsize
\scalebox{0.95}{
\begin{tabular}{>{\RaggedRight}p{4cm}rrr>{\RaggedRight}p{4cm}}\toprule
Null hypothesis &Instances &Test &p-value &Decision \\\cmidrule{1-5}
The mean for each population is equal &x7day\_small &Kruskal-Wallis &0.0000* &We reject $\text{H}_o$ and accept $\text{H}_a$: at least one population mean different from the rest \\
The mean for each population is equal &x7day\_medium &Kruskal-Wallis &0.0000* &We reject $\text{H}_o$ and accept $\text{H}_a$: at least one population mean different from the rest \\
The mean for each population is equal &x7day\_large &Kruskal-Wallis &0.0000* &We reject $\text{H}_o$ and accept $\text{H}_a$: at least one population mean different from the rest \\
\bottomrule
*  Significance level 0.05
\end{tabular}
}

\end{table}
\begin{table}[H]\centering
\caption{Differences mean test summary between the constructive methods HARDUC and HGPS.}\label{tabla:Differences_constructive}
\scriptsize
\scalebox{0.95}{
\begin{tabular}{>{\RaggedRight}p{4.3cm}rrr>{\RaggedRight}p{4cm}}\toprule
Null hypothesis &Instances&Test &p-value &Decision \\\cmidrule{1-5}
The means difference of the samples from the same distribution &x7day\_small &Mann-Whitney &0.0002* &We reject $\text{H}_o$ and accept $\text{H}_a$: HARDUC's mean is less than HGPS's mean \\
The means difference of the samples from the same distribution &x7day\_medium &Mann-Whitney &0.0000* &We reject $\text{H}_o$ and accept $\text{H}_a$: HARDUC's mean is less than HGPS's mean \\
The means difference of the samples from the same distribution &x7day\_large &Mann-Whitney &0.0000* &We reject $\text{H}_o$ and accept $\text{H}_a$: HARDUC's mean is less than HGPS's mean \\
\bottomrule
*  Significance level 0.05
\end{tabular}
}
\end{table}

%%%%%%%%%%%%%%%%%%%%%%%%%%%%%%%%%%%%%%%%%%%%%%%%%%%%%%%
%%%%%%%%%%%%%%%%%%%%%%%%%%%%%%%%%%%%%%%%%%%%%%%%%%%%%%%
\setcounter{table}{0}
\setcounter{figure}{0}
\renewcommand{\thetable}{B\arabic{table}}
\renewcommand{\thefigure}{B\arabic{figure}}
%\subsection*{Appendix B}\label{app:appB}
\section{Results for Local Branching and Kernel Search}\label{app:appB}


We perform statistical tests by analyzing variance and comparing means to determine whether the improvements are significant. The summary of the results of the analysis of variance for the three groups of instances can be found in Table \ref{tabla:ANOVA_allmethods_1h}.

%%%%%%%%%%%%%%%%%%%%%%%%%%%%%%%%%%%%%%%%%%%%%%%%%%%%%%
\begin{table}[!htp]\centering
\caption{Analysis of variance summary for LB1\_1h, LB2\_1h, LB3\_1h, LB4\_1h, KS\_1h, SM1\_1h, SM2\_1h under a running time limit of 4000 seconds.}\label{tabla:ANOVA_allmethods_1h}
\scriptsize
\scalebox{0.95}{
\begin{tabular}{>{\RaggedRight}p{4cm}rrr>{\RaggedRight}p{4cm}}\toprule
Null hypothesis &Instances &Test &p-value &Decision \\\cmidrule{1-5}
The mean for each population is equal &x7day\_small &Kruskal-Wallis &0.9758 &We fail to reject $\text{H}_o$ \\
The mean for each population is equal &x7day\_medium &Kruskal-Wallis &0.0000* &We reject $\text{H}_o$ and accept $\text{H}_a$: at least one population mean different from the rest \\
The mean for each population is equal &x7day\_large &Kruskal-Wallis &0.0066* &We reject $\text{H}_o$ and accept $\text{H}_a$: at least one population mean different from the rest \\
\bottomrule
*  Significance level 0.05
\end{tabular}
}
\end{table}

The analysis of variance results rejects the null hypothesis of equality of the means and does not reject the alternative hypothesis that at least one of the means differs from that of the rest of the methods.

The results of the mean comparison study among all the methods, for 4000 seconds of execution and for the three groups of instances are presented in Tables \ref{tabla:Differences_small_1h}, \ref{tabla:Differences_medium_1h}, and \ref{tabla:Differences_large_1h}.

%%%%%

We carry out statistical testing to determine whether the improvements are significant by analyzing variance and comparing means. The summary of the results of the analysis of variance for the three groups of instances can be found in Table \ref{tabla:ANOVA_allmethods_2h}.

\begin{table}[!htp]\centering
\caption{Analysis of variance summary for LB1, LB2, LB3, LB4, KS, MILP, SM2 under a running time limit of 7200 seconds.}\label{tabla:ANOVA_allmethods_2h}
\scriptsize
\scalebox{0.95}{
\begin{tabular}{>{\RaggedRight}p{5cm}rrr>{\RaggedRight}p{4cm}r}\toprule
Null hypothesis &Instances &Test &p-value &Decision \\\cmidrule{1-5}
The mean for each population is equal &x7day\_small &Kruskal-Wallis &0.7289 &We fail to reject $\text{H}_o$ \\
The mean for each population is equal &x7day\_medium &Kruskal-Wallis &0.0000* &We reject $\text{H}_o$ and accept $\text{H}_a$: at least one population mean different from the rest \\
The mean for each population is equal &x7day\_large &Kruskal-Wallis &0.0007* &We reject $\text{H}_o$ and accept $\text{H}_a$: at least one population mean different from the rest \\
\bottomrule
*  Significance level 0.05
\end{tabular}
}
\end{table}


%%%%%%%%%%%%%%%%%%%%%%%%%%%%%%%%%%%%%%%%%%%%%%%%%%%%%%%
%%%%%%%%%%%%%%%%%%%%%%%%%%%%%%%%%%%%%%%%%%%%%%%%%%%%%%%
\setcounter{table}{0}
\setcounter{figure}{0}
\renewcommand{\thetable}{C\arabic{table}}
\renewcommand{\thefigure}{C\arabic{figure}}
%\subsection*{Appendix C} \label{app:appC}
\section{Statistical Tests} \label{app:appC}

\begin{table}[H]\centering
\caption{Differences mean test summary between the constructive methods HARDUC and HGPS.}\label{tabla:Differences_constructive2}
\scriptsize
\scalebox{0.95}{
\begin{tabular}{>{\RaggedRight}p{4.3cm}rrr>{\RaggedRight}p{4cm}}\toprule
Null hypothesis &Instances&Test &p-value &Decision \\\cmidrule{1-5}
The means difference of the samples from the same distribution &x7day\_small &Mann-Whitney &0.0002* &We reject $\text{H}_o$ and accept $\text{H}_a$: HARDUC's mean is less than HGPS's mean \\
The means difference of the samples from the same distribution &x7day\_medium &Mann-Whitney &0.0000* &We reject $\text{H}_o$ and accept $\text{H}_a$: HARDUC's mean is less than HGPS's mean \\
The means difference of the samples from the same distribution &x7day\_large &Mann-Whitney &0.0000* &We reject $\text{H}_o$ and accept $\text{H}_a$: HARDUC's mean is less than HGPS's mean \\
\bottomrule
*  Significance level 0.05
\end{tabular}
}
\end{table}


%%%%%%%%%%%%%%%%%%%%%%%%%%%%%%%%%%%%%%%%%%%%%%%%%%%%%%
%Means difference hypothesis test summary (x7day\_small) 1 hora
\begin{table}[H]\centering
\caption{Means difference statistical test summary among all methods for instances from group x7day\_small under a running time limit of 4000 seconds.}
\label{tabla:Differences_small_1h}
\scriptsize
\scalebox{0.82}{
\begin{tabular}{lrr>{\RaggedRight}p{2.7cm}}\toprule
Null hypothesis &Test &p-value &Decision \\\cmidrule{1-4}
SM1\_1h-LB2\_1h: There is no difference between the two population means & Mann-Whitney & 0.3779 & We fail to reject $\text{H}_o$ \\
SM1\_1h-KS\_1h: There is no difference between the two population means & T-test for two samples & 0.2725 & We fail to reject $\text{H}_o$ \\
SM1\_1h-LB1\_1h: There is no difference between the two population means & T-test for two samples & 0.2499 & We fail to reject $\text{H}_o$ \\
SM1\_1h-LB4\_1h: There is no difference between the two population means & T-test for two samples & 0.1737 & We fail to reject $\text{H}_o$ \\
SM1\_1h-LB3\_1h: There is no difference between the two population means & T-test for two samples & 0.1318 & We fail to reject $\text{H}_o$ \\
SM1\_1h-SM2\_1h: There is no difference between the two population means & T-test for two samples & 0.0888 & We fail to reject $\text{H}_o$ \\
LB2\_1h-KS\_1h: There is no difference between the two population means & Mann-Whitney & 0.5484 & We fail to reject $\text{H}_o$ \\
LB2\_1h-LB1\_1h: There is no difference between the two population means & Mann-Whitney & 0.4302 & We fail to reject $\text{H}_o$ \\
LB2\_1h-LB4\_1h: There is no difference between the two population means & Mann-Whitney & 0.2714 & We fail to reject $\text{H}_o$ \\
LB2\_1h-LB3\_1h: There is no difference between the two population means & Mann-Whitney & 0.2367 & We fail to reject $\text{H}_o$ \\
LB2\_1h-SM2\_1h: There is no difference between the two population means & Mann-Whitney & 0.1617 & We fail to reject $\text{H}_o$ \\
KS\_1h-LB1\_1h: There is no difference between the two population means & T-test for two samples & 0.492 & We fail to reject $\text{H}_o$ \\
KS\_1h-LB4\_1h: There is no difference between the two population means & T-test for two samples & 0.4016 & We fail to reject $\text{H}_o$ \\
KS\_1h-LB3\_1h: There is no difference between the two population means & T-test for two samples & 0.3336 & We fail to reject $\text{H}_o$ \\
KS\_1h-SM2\_1h: There is no difference between the two population means & T-test for two samples & 0.2923 & We fail to reject $\text{H}_o$ \\
LB1\_1h-LB4\_1h: There is no difference between the two population means & T-test for two samples & 0.4039 & We fail to reject $\text{H}_o$ \\
LB1\_1h-LB3\_1h: There is no difference between the two population means & T-test for two samples & 0.332 & We fail to reject $\text{H}_o$ \\
LB1\_1h-SM2\_1h: There is no difference between the two population means & T-test for two samples & 0.287 & We fail to reject $\text{H}_o$ \\
LB4\_1h-LB3\_1h: There is no difference between the two population means & T-test for two samples & 0.4211 & We fail to reject $\text{H}_o$ \\
LB4\_1h-SM2\_1h: There is no difference between the two population means & T-test for two samples & 0.3808 & We fail to reject $\text{H}_o$ \\
LB3\_1h-SM2\_1h: There is no difference between the two population means & T-test for two samples & 0.4675 & We fail to reject $\text{H}_o$ \\
\bottomrule
*  Significance level 0.05
\end{tabular}
}
\end{table}



%Means difference hypothesis test summary (x7day\_small) 2 hours
\begin{table}[H]\centering
\caption{Means difference statistical test summary among all methods for instances from group x7day\_small under a running time limit of 7200 seconds.}
\label{tabla:Differences_small_2h}
\scriptsize
\scalebox{0.9}{
\begin{tabular}{>{\RaggedRight}p{5.5cm}rr>{\RaggedRight}p{3.5cm}}
\toprule
Null hypothesis &Test &p-value &Decision \\\cmidrule{1-4}
SM1-LB2: There is no difference between the two population means & Mann-Whitney & 0.1584 & We fail to reject $\text{H}_o$ \\
SM1-LB1: There is no difference between the two population means & Mann-Whitney & 0.1396 & We fail to reject $\text{H}_o$ \\
SM1-LB4: There is no difference between the two population means & Mann-Whitney & 0.117 & We fail to reject $\text{H}_o$ \\
SM1-SM2: There is no difference between the two population means & T-test for two samples & *0.0362 & We reject $\text{H}_o$ and accept $\text{H}_a$: MILP's mean less than MILP2's \\
SM1-KS: There is no difference between the two population means & T-test for two samples & 0.0501 & We fail to reject $\text{H}_o$ \\
SM1-LB3: There is no difference between the two population means & T-test for two samples & *0.0358 & We reject $\text{H}_o$ and accept $\text{H}_a$: MILP's mean less than LB3's \\
LB2-LB1: There is no difference between the two population means & Mann-Whitney & 0.4569 & We fail to reject $\text{H}_o$ \\
LB2-LB4: There is no difference between the two population means & Mann-Whitney & 0.4143 & We fail to reject $\text{H}_o$ \\
LB2-SM2: There is no difference between the two population means & Mann-Whitney & 0.2326 & We fail to reject $\text{H}_o$ \\
LB2-KS: There is no difference between the two population means & Mann-Whitney & 0.4091 & We fail to reject $\text{H}_o$ \\
LB2-LB3: There is no difference between the two population means & Mann-Whitney & 0.2989 & We fail to reject $\text{H}_o$ \\
LB1-LB4: There is no difference between the two population means & Mann-Whitney & 0.4515 & We fail to reject $\text{H}_o$ \\
LB1-SM2: There is no difference between the two population means & Mann-Whitney & 0.2669 & We fail to reject $\text{H}_o$ \\
LB1-KS: There is no difference between the two population means & Mann-Whitney & 0.4623 & We fail to reject $\text{H}_o$ \\
LB1-LB3: There is no difference between the two population means & Mann-Whitney & 0.2896 & We fail to reject $\text{H}_o$ \\
LB4-SM2: There is no difference between the two population means & Mann-Whitney & 0.3036 & We fail to reject $\text{H}_o$ \\
LB4-KS: There is no difference between the two population means & Mann-Whitney & 0.4623 & We fail to reject $\text{H}_o$ \\
LB4-LB3: There is no difference between the two population means & Mann-Whitney & 0.3375 & We fail to reject $\text{H}_o$ \\
SM2-KS: There is no difference between the two population means & T-test for two samples & 0.3651 & We fail to reject $\text{H}_o$ \\
SM2-LB3: There is no difference between the two population means & T-test for two samples & 0.3434 & We fail to reject $\text{H}_o$ \\
KS-LB3: There is no difference between the two population means & T-test for two samples & 0.4909 & We fail to reject $\text{H}_o$ \\
\bottomrule
*  Significance level 0.05
\end{tabular}
}
\end{table}

%%%%%%%%%%%%%%%%%%%%%%%%%%%%%%%%%%%%%%%%%%%%%%%%%%%%%%
%Means difference hypothesis test summary (x7day\_medium)
\begin{table}[H]\centering
\caption{Means difference statistical test summary among all methods for instances from group x7day\_medium under a running time limit of 4000 seconds.}
\label{tabla:Differences_medium_1h}
\scriptsize
\scalebox{0.85}{
\begin{tabular}{>{\RaggedRight}p{5.5cm}rr>{\RaggedRight}p{5.5cm}r}\toprule
Null hypothesis &Test &p-value &Decision \\\cmidrule{1-4}
KS\_1h-LB2\_1h: There is no difference between the two population means & Mann-Whitney & *0.0108 & We reject $\text{H}_o$ and accept $\text{H}_a$: KS\_1h's mean less than LB2\_1h's \\
KS\_1h-LB1\_1h: There is no difference between the two population means & Mann-Whitney & *0.0031 & We reject $\text{H}_o$ and accept $\text{H}_a$: KS\_1h's mean less than LB1\_1h's \\
KS\_1h-LB4\_1h: There is no difference between the two population means & Mann-Whitney & *0.0022 & We reject $\text{H}_o$ and accept $\text{H}_a$: KS\_1h's mean less than LB4\_1h's \\
KS\_1h-LB3\_1h: There is no difference between the two population means & Mann-Whitney & *0.001 & We reject $\text{H}_o$ and accept $\text{H}_a$: KS\_1h's mean less than LB3\_1h's \\
KS\_1h-SM1\_1h: There is no difference between the two population means & Mann-Whitney & *0.0 & We reject $\text{H}_o$ and accept $\text{H}_a$: KS\_1h's mean less than SM1\_1h's \\
KS\_1h-SM2\_1h: There is no difference between the two population means & Mann-Whitney & *0.0 & We reject $\text{H}_o$ and accept $\text{H}_a$: KS\_1h's mean less than SM2\_1h's \\
LB2\_1h-LB1\_1h: There is no difference between the two population means & Mann-Whitney & 0.2699 & We fail to reject $\text{H}_o$ \\
LB2\_1h-LB4\_1h: There is no difference between the two population means & Mann-Whitney & 0.2405 & We fail to reject $\text{H}_o$ \\
LB2\_1h-LB3\_1h: There is no difference between the two population means & Mann-Whitney & 0.1593 & We fail to reject $\text{H}_o$ \\
LB2\_1h-SM1\_1h: There is no difference between the two population means & Mann-Whitney & *0.0002 & We reject $\text{H}_o$ and accept $\text{H}_a$: LB2\_1h's mean less than SM1\_1h's \\
LB2\_1h-SM2\_1h: There is no difference between the two population means & Mann-Whitney & *0.0002 & We reject $\text{H}_o$ and accept $\text{H}_a$: LB2\_1h's mean less than SM2\_1h's \\
LB1\_1h-LB4\_1h: There is no difference between the two population means & Mann-Whitney & 0.4745 & We fail to reject $\text{H}_o$ \\
LB1\_1h-LB3\_1h: There is no difference between the two population means & Mann-Whitney & 0.3538 & We fail to reject $\text{H}_o$ \\
LB1\_1h-SM1\_1h: There is no difference between the two population means & Mann-Whitney & *0.0003 & We reject $\text{H}_o$ and accept $\text{H}_a$: LB1\_1h's mean less than SM1\_1h's \\
LB1\_1h-SM2\_1h: There is no difference between the two population means & Mann-Whitney & *0.0003 & We reject $\text{H}_o$ and accept $\text{H}_a$: LB1\_1h's mean less than SM2\_1h's \\
LB4\_1h-LB3\_1h: There is no difference between the two population means & Mann-Whitney & 0.3606 & We fail to reject $\text{H}_o$ \\
LB4\_1h-SM1\_1h: There is no difference between the two population means & Mann-Whitney & *0.0005 & We reject $\text{H}_o$ and accept $\text{H}_a$: LB4\_1h's mean less than SM1\_1h's \\
LB4\_1h-SM2\_1h: There is no difference between the two population means & Mann-Whitney & *0.0004 & We reject $\text{H}_o$ and accept $\text{H}_a$: LB4\_1h's mean less than SM2\_1h's \\
LB3\_1h-SM1\_1h: There is no difference between the two population means & Mann-Whitney & *0.0009 & We reject $\text{H}_o$ and accept $\text{H}_a$: LB3\_1h's mean less than SM1\_1h's \\
LB3\_1h-SM2\_1h: There is no difference between the two population means & Mann-Whitney & *0.0005 & We reject $\text{H}_o$ and accept $\text{H}_a$: LB3\_1h's mean less than SM2\_1h's \\
SM1\_1h-SM2\_1h: There is no difference between the two population means & Mann-Whitney & 0.2699 & We fail to reject $\text{H}_o$ \\
\bottomrule
*  Significance level 0.05
\end{tabular}
}
\end{table}


%%%%%%%%%%%%%%%%%%%%%%%%%%%%%%%%%%%%%%%%%%%%%%%%%
%Means difference hypothesis test summary (x7day\_medium)
\begin{table}[H]\centering
\caption{Means difference statistical test summary among all methods for instances from group x7day\_medium under a running time limit of 7200 seconds.}
\label{tabla:Differences_medium_2h}
\scriptsize
\scalebox{0.85}{
\begin{tabular}{>{\RaggedRight}p{5.5cm}rr>{\RaggedRight}p{5.5cm}r}\toprule
Null hypothesis &Test &p-value &Decision \\\cmidrule{1-4}
KS-SM1: There is no difference between the two population means & Mann-Whitney & 0.3813 & We fail to reject $\text{H}_o$ \\
KS-LB1: There is no difference between the two population means & Mann-Whitney & *0.0165 & We reject $\text{H}_o$ and accept $\text{H}_a$: KS's mean less than LB1's \\
KS-LB4: There is no difference between the two population means & Mann-Whitney & *0.0125 & We reject $\text{H}_o$ and accept $\text{H}_a$: KS's mean less than LB4's \\
KS-LB2: There is no difference between the two population means & Mann-Whitney & *0.0131 & We reject $\text{H}_o$ and accept $\text{H}_a$: KS's mean less than LB2's \\
KS-LB3: There is no difference between the two population means & Mann-Whitney & *0.0012 & We reject $\text{H}_o$ and accept $\text{H}_a$: KS's mean less than LB3's \\
KS-SM2: There is no difference between the two population means & Mann-Whitney & *0.0 & We reject $\text{H}_o$ and accept $\text{H}_a$: KS's mean less than MILP2's \\
SM1-LB1: There is no difference between the two population means & Mann-Whitney & 0.061 & We fail to reject $\text{H}_o$ \\
SM1-LB4: There is no difference between the two population means & Mann-Whitney & 0.061 & We fail to reject $\text{H}_o$ \\
SM1-LB2: There is no difference between the two population means & Mann-Whitney & 0.0691 & We fail to reject $\text{H}_o$ \\
SM1-LB3: There is no difference between the two population means & Mann-Whitney & *0.0103 & We reject $\text{H}_o$ and accept $\text{H}_a$: MILP's mean less than LB3's \\
SM1-SM2: There is no difference between the two population means & Mann-Whitney & *0.0 & We reject $\text{H}_o$ and accept $\text{H}_a$: MILP's mean less than MILP2's \\
LB1-LB4: There is no difference between the two population means & Mann-Whitney & 0.4418 & We fail to reject $\text{H}_o$ \\
LB1-LB2: There is no difference between the two population means & Mann-Whitney & 0.5219 & We fail to reject $\text{H}_o$ \\
LB1-LB3: There is no difference between the two population means & Mann-Whitney & 0.1463 & We fail to reject $\text{H}_o$ \\
LB1-SM2: There is no difference between the two population means & Mann-Whitney & *0.0001 & We reject $\text{H}_o$ and accept $\text{H}_a$: LB1's mean less than MILP2's \\
LB4-LB2: There is no difference between the two population means & Mann-Whitney & 0.551 & We fail to reject $\text{H}_o$ \\
LB4-LB3: There is no difference between the two population means & T-test for two samples & 0.2583 & We fail to reject $\text{H}_o$ \\
LB4-SM2: There is no difference between the two population means & T-test for two samples & *0.0001 & We reject $\text{H}_o$ and accept $\text{H}_a$: LB4's mean less than MILP2's \\
LB2-LB3: There is no difference between the two population means & Mann-Whitney & 0.1638 & We fail to reject $\text{H}_o$ \\
LB2-SM2: There is no difference between the two population means & Mann-Whitney & *0.0002 & We reject $\text{H}_o$ and accept $\text{H}_a$: LB2's mean less than MILP2's \\
LB3-SM2: There is no difference between the two population means & Mann-Whitney & *0.0004 & We reject $\text{H}_o$ and accept $\text{H}_a$: LB3's mean less than MILP2's \\
\bottomrule
*  Significance level 0.05
\end{tabular}
}
\end{table}



%%%%%%%%%%%%%%%%%%%%%%%%%%%%%%%%%%%%%%%%%%%%%
%Means difference hypothesis test summary (x7day\_large)
\begin{table}[H]\centering
\caption{Means difference statistical test summary among all methods for instances from group x7day\_large under a running time limit of 4000 seconds.}
\label{tabla:Differences_large_1h}
\scriptsize
\scalebox{0.85}{
\begin{tabular}{>{\RaggedRight}p{5.5cm}rr>{\RaggedRight}p{5.5cm}r}\toprule
Null hypothesis &Test &p-value &Decision \\\cmidrule{1-4}
KS\_1h-LB1\_1h: There is no difference between the two population means & T-test for two samples & *0.0071 & We reject $\text{H}_o$ and accept $\text{H}_a$: KS\_1h's mean less than LB1\_1h's \\
KS\_1h-LB2\_1h: There is no difference between the two population means & T-test for two samples & *0.0038 & We reject $\text{H}_o$ and accept $\text{H}_a$: KS\_1h's mean less than LB2\_1h's \\
KS\_1h-LB4\_1h: There is no difference between the two population means & T-test for two samples & *0.0005 & We reject $\text{H}_o$ and accept $\text{H}_a$: KS\_1h's mean less than LB4\_1h's \\
KS\_1h-LB3\_1h: There is no difference between the two population means & Mann-Whitney & *0.0021 & We reject $\text{H}_o$ and accept $\text{H}_a$: KS\_1h's mean less than LB3\_1h's \\
KS\_1h-SM1\_1h: There is no difference between the two population means & Mann-Whitney & *0.0003 & We reject $\text{H}_o$ and accept $\text{H}_a$: KS\_1h's mean less than SM1\_1h's \\
KS\_1h-SM2\_1h: There is no difference between the two population means & Mann-Whitney & *0.0001 & We reject $\text{H}_o$ and accept $\text{H}_a$: KS\_1h's mean less than SM2\_1h's \\
LB1\_1h-LB2\_1h: There is no difference between the two population means & T-test for two samples & 0.2958 & We fail to reject $\text{H}_o$ \\
LB1\_1h-LB4\_1h: There is no difference between the two population means & T-test for two samples & 0.1441 & We fail to reject $\text{H}_o$ \\
LB1\_1h-LB3\_1h: There is no difference between the two population means & T-test for two samples & 0.1603 & We fail to reject $\text{H}_o$ \\
LB1\_1h-SM1\_1h: There is no difference between the two population means & Mann-Whitney & *0.0425 & We reject $\text{H}_o$ and accept $\text{H}_a$: LB1\_1h's mean less than SM1\_1h's \\
LB1\_1h-SM2\_1h: There is no difference between the two population means & T-test for two samples & *0.0033 & We reject $\text{H}_o$ and accept $\text{H}_a$: LB1\_1h's mean less than SM2\_1h's \\
LB2\_1h-LB4\_1h: There is no difference between the two population means & T-test for two samples & 0.3234 & We fail to reject $\text{H}_o$ \\
LB2\_1h-LB3\_1h: There is no difference between the two population means & T-test for two samples & 0.329 & We fail to reject $\text{H}_o$ \\
LB2\_1h-SM1\_1h: There is no difference between the two population means & Mann-Whitney & 0.1208 & We fail to reject $\text{H}_o$ \\
LB2\_1h-SM2\_1h: There is no difference between the two population means & T-test for two samples & *0.0155 & We reject $\text{H}_o$ and accept $\text{H}_a$: LB2\_1h's mean less than SM2\_1h's \\
LB4\_1h-LB3\_1h: There is no difference between the two population means & T-test for two samples & 0.4934 & We fail to reject $\text{H}_o$ \\
LB4\_1h-SM1\_1h: There is no difference between the two population means & Mann-Whitney & 0.2347 & We fail to reject $\text{H}_o$ \\
LB4\_1h-SM2\_1h: There is no difference between the two population means & T-test for two samples & *0.032 & We reject $\text{H}_o$ and accept $\text{H}_a$: LB4\_1h's mean less than SM2\_1h's \\
LB3\_1h-SM1\_1h: There is no difference between the two population means & Mann-Whitney & 0.1853 & We fail to reject $\text{H}_o$ \\
LB3\_1h-SM2\_1h: There is no difference between the two population means & T-test for two samples & *0.0408 & We reject $\text{H}_o$ and accept $\text{H}_a$: LB3\_1h's mean less than SM2\_1h's \\
SM1\_1h-SM2\_1h: There is no difference between the two population means & Mann-Whitney & 0.1208 & We fail to reject $\text{H}_o$ \\
\bottomrule
*  Significance level 0.05
\end{tabular}
}
\end{table}


%%%%%%%%%%%%%%%%%%%%%%%%%%%%%%%%%%%%%%%%%%%%%%%%%%%%%
%Means difference hypothesis test summary (x7day\_large)
\begin{table}[H]\centering
\caption{Means difference statistical test summary among all methods for instances from group x7day\_large under a running time limit of 7200 seconds.}
\label{tabla:Differences_large_2h}
\scriptsize
\scalebox{0.85}{
\begin{tabular}{>{\RaggedRight}p{5.5cm}rr>{\RaggedRight}p{5.5cm}r}\toprule
Null hypothesis &Test &p-value &Decision \\\cmidrule{1-4}
KS-LB2: There is no difference between the two population means & Mann-Whitney & *0.0164 & We reject $\text{H}_o$ and accept $\text{H}_a$: KS's mean less than LB2's \\
KS-LB1: There is no difference between the two population means & Mann-Whitney & *0.0105 & We reject $\text{H}_o$ and accept $\text{H}_a$: KS's mean less than LB1's \\
KS-LB4: There is no difference between the two population means & Mann-Whitney & *0.0072 & We reject $\text{H}_o$ and accept $\text{H}_a$: KS's mean less than LB4's \\
KS-LB3: There is no difference between the two population means & Mann-Whitney & *0.0019 & We reject $\text{H}_o$ and accept $\text{H}_a$: KS's mean less than LB3's \\
KS-SM1: There is no difference between the two population means & Mann-Whitney & *0.0021 & We reject $\text{H}_o$ and accept $\text{H}_a$: KS's mean less than MILP's \\
KS-SM2: There is no difference between the two population means & Mann-Whitney & *0.0001 & We reject $\text{H}_o$ and accept $\text{H}_a$: KS's mean less than MILP2's \\
LB2-LB1: There is no difference between the two population means & T-test for two samples & 0.3812 & We fail to reject $\text{H}_o$ \\
LB2-LB4: There is no difference between the two population means & T-test for two samples & 0.2823 & We fail to reject $\text{H}_o$ \\
LB2-LB3: There is no difference between the two population means & T-test for two samples & 0.1417 & We fail to reject $\text{H}_o$ \\
LB2-SM1: There is no difference between the two population means & T-test for two samples & 0.0858 & We fail to reject $\text{H}_o$ \\
LB2-SM2: There is no difference between the two population means & T-test for two samples & *0.0012 & We reject $\text{H}_o$ and accept $\text{H}_a$: LB2's mean less than MILP2's \\
LB1-LB4: There is no difference between the two population means & T-test for two samples & 0.3952 & We fail to reject $\text{H}_o$ \\
LB1-LB3: There is no difference between the two population means & T-test for two samples & 0.2175 & We fail to reject $\text{H}_o$ \\
LB1-SM1: There is no difference between the two population means & T-test for two samples & 0.1337 & We fail to reject $\text{H}_o$ \\
LB1-SM2: There is no difference between the two population means & T-test for two samples & *0.0022 & We reject $\text{H}_o$ and accept $\text{H}_a$: LB1's mean less than MILP2's \\
LB4-LB3: There is no difference between the two population means & T-test for two samples & 0.2932 & We fail to reject $\text{H}_o$ \\
LB4-SM1: There is no difference between the two population means & T-test for two samples & 0.1826 & We fail to reject $\text{H}_o$ \\
LB4-SM2: There is no difference between the two population means & T-test for two samples & *0.0032 & We reject $\text{H}_o$ and accept $\text{H}_a$: LB4's mean less than MILP2's \\
LB3-SM1: There is no difference between the two population means & T-test for two samples & 0.3414 & We fail to reject $\text{H}_o$ \\
LB3-SM2: There is no difference between the two population means & T-test for two samples & *0.0101 & We reject $\text{H}_o$ and accept $\text{H}_a$: LB3's mean less than MILP2's \\
SM1-SM2: There is no difference between the two population means & T-test for two samples & *0.0290 & We reject $\text{H}_o$ and accept $\text{H}_a$: MILP's mean less than MILP2's \\
\bottomrule
*  Significance level 0.05
\end{tabular}
}
\end{table}

%%%=========================================================App D
\newpage
\setcounter{table}{0}
\setcounter{figure}{0}
\renewcommand{\thetable}{D\arabic{table}}
\renewcommand{\thefigure}{D\arabic{figure}}
\section{Description of instances} \label{app:appD}


\begin{table}[!htp]\centering
\caption{x7day\_small instances group.}\label{tab:firstinst}
\scriptsize
\scalebox{0.9}{
\begin{tabular}{lrrrrrrrrrr}\toprule
Instance & Gen & \multicolumn{8}{c}{Configuration} \\\cmidrule{3-10}
 & & G1 & G2 & G3 & G4 & G5 & G6 & G7 & G8 \\\cmidrule{1-10}
061 & 28 & 12 & 11 & 0 & 0 & 1 & 4 & 0 & 0 \\
062 & 35 & 13 & 15 & 2 & 0 & 4 & 0 & 0 & 1 \\
063 & 44 & 15 & 13 & 2 & 6 & 3 & 1 & 1 & 3 \\
064 & 45 & 15 & 11 & 0 & 1 & 4 & 5 & 6 & 3 \\
065 & 49 & 15 & 13 & 3 & 7 & 5 & 3 & 2 & 1 \\
066 & 50 & 10 & 10 & 2 & 5 & 7 & 5 & 6 & 5 \\
067 & 51 & 17 & 16 & 1 & 3 & 1 & 7 & 2 & 4 \\
068 & 51 & 17 & 10 & 6 & 5 & 2 & 1 & 3 & 7 \\
069 & 52 & 12 & 17 & 4 & 7 & 5 & 2 & 0 & 5 \\
070 & 54 & 13 & 12 & 5 & 7 & 2 & 5 & 4 & 6 \\
101 & 60 & 11 & 18 & 6 & 3 & 7 & 8 & 2 & 5 \\
102 & 62 & 16 & 17 & 5 & 3 & 5 & 6 & 7 & 3 \\
103 & 64 & 13 & 20 & 4 & 9 & 4 & 6 & 2 & 6 \\
104 & 66 & 17 & 16 & 8 & 8 & 0 & 6 & 6 & 5 \\
105 & 70 & 14 & 21 & 9 & 5 & 8 & 6 & 3 & 4 \\
106 & 71 & 25 & 18 & 2 & 4 & 7 & 5 & 2 & 8 \\
107 & 75 & 16 & 25 & 5 & 9 & 3 & 3 & 8 & 6 \\
108 & 76 & 21 & 21 & 4 & 3 & 8 & 10 & 4 & 5 \\
109 & 80 & 17 & 24 & 8 & 9 & 8 & 1 & 8 & 5 \\
110 & 81 & 22 & 21 & 4 & 9 & 4 & 3 & 8 & 10 \\
\bottomrule
\end{tabular}
}
\end{table}

\begin{table}[!htp]\centering
\caption{x7day\_medium instances group.}\label{tab:secondinst}
\scriptsize
\scalebox{0.8}{
\begin{tabular}{lrrrrrrrrrr}\toprule
Instance & Gen & \multicolumn{8}{c}{Configuration} \\\cmidrule{3-10}
 & & G1 & G2 & G3 & G4 & G5 & G6 & G7 & G8 \\\cmidrule{1-10}
111 & 85 & 19 & 27 & 6 & 6 & 8 & 9 & 3 & 7 \\
112 & 86 & 24 & 23 & 6 & 7 & 4 & 8 & 6 & 8 \\
113 & 90 & 24 & 28 & 6 & 7 & 6 & 3 & 12 & 4 \\
114 & 91 & 26 & 24 & 9 & 4 & 5 & 10 & 6 & 7 \\
115* & 94 & 22 & 28 & 4 & 10 & 5 & 9 & 7 & 9 \\
116 & 97 & 27 & 26 & 8 & 14 & 5 & 6 & 7 & 4 \\
117 & 98 & 24 & 30 & 4 & 8 & 10 & 7 & 10 & 5 \\
118* & 102 & 29 & 28 & 10 & 6 & 5 & 9 & 9 & 6 \\
119* & 104 & 25 & 32 & 9 & 0 & 10 & 11 & 9 & 8 \\
120* & 106 & 30 & 29 & 11 & 9 & 5 & 2 & 11 & 9 \\
121* & 109 & 32 & 34 & 9 & 7 & 8 & 9 & 4 & 6 \\
122 & 111 & 32 & 36 & 6 & 4 & 10 & 7 & 10 & 6 \\
123 & 114 & 29 & 36 & 11 & 10 & 6 & 5 & 7 & 10 \\
124 & 116 & 34 & 32 & 7 & 5 & 9 & 9 & 11 & 9 \\
125 & 119 & 32 & 37 & 9 & 5 & 7 & 11 & 9 & 9 \\
126 & 121 & 35 & 34 & 7 & 11 & 12 & 4 & 7 & 11 \\
127* & 125 & 32 & 39 & 5 & 8 & 5 & 15 & 13 & 8 \\
128* & 126 & 37 & 36 & 10 & 12 & 4 & 12 & 10 & 5 \\
129 & 129 & 33 & 40 & 9 & 7 & 7 & 14 & 5 & 14 \\
130 & 131 & 38 & 37 & 14 & 10 & 6 & 7 & 5 & 14 \\
131 & 131 & 34 & 42 & 8 & 8 & 8 & 11 & 14 & 6 \\
071 & 132 & 46 & 45 & 8 & 0 & 5 & 0 & 12 & 16 \\
132 & 135 & 37 & 44 & 8 & 8 & 9 & 9 & 13 & 7 \\
133 & 138 & 42 & 40 & 7 & 11 & 9 & 10 & 6 & 13 \\
134 & 140 & 38 & 45 & 8 & 7 & 10 & 14 & 10 & 8 \\
135 & 143 & 43 & 42 & 7 & 12 & 7 & 15 & 10 & 7 \\
136 & 146 & 40 & 47 & 9 & 7 & 15 & 13 & 7 & 8 \\
137 & 148 & 45 & 44 & 7 & 16 & 9 & 10 & 10 & 7 \\
138 & 151 & 41 & 48 & 9 & 16 & 11 & 8 & 9 & 9 \\
139 & 154 & 46 & 45 & 10 & 15 & 10 & 3 & 15 & 10 \\
140 & 156 & 43 & 48 & 16 & 11 & 9 & 12 & 9 & 8 \\
072 & 156 & 40 & 54 & 14 & 8 & 3 & 15 & 9 & 13 \\
073 & 156 & 50 & 41 & 19 & 11 & 4 & 4 & 12 & 15 \\
\bottomrule
\end{tabular}
}
\scriptsize\caption*{*Solver did not produce an initial integer solution}
\end{table}

\begin{table}[!htp]\centering
\caption{x7day\_large instances group.}\label{tab:thirdinst}
\scriptsize
\scalebox{0.9}{
\begin{tabular}{lrrrrrrrrrr}\toprule
Instance & Gen & \multicolumn{8}{c}{Configuration} \\\cmidrule{3-10}
 & & G1 & G2 & G3 & G4 & G5 & G6 & G7 & G8 \\\cmidrule{1-10}
074 & 165 & 51 & 58 & 17 & 19 & 16 & 1 & 2 & 1 \\
075 & 167 & 43 & 46 & 17 & 15 & 13 & 15 & 6 & 12 \\
076 & 172 & 50 & 59 & 8 & 15 & 1 & 18 & 4 & 17 \\
077 & 182 & 53 & 50 & 17 & 15 & 16 & 5 & 14 & 12 \\
078 & 182 & 45 & 57 & 19 & 7 & 19 & 19 & 5 & 11 \\
079 & 183 & 58 & 50 & 15 & 7 & 16 & 18 & 7 & 12 \\
080 & 187 & 55 & 48 & 18 & 5 & 18 & 17 & 15 & 11 \\
141* & 185 & 52 & 57 & 15 & 12 & 12 & 13 & 11 & 13 \\
142* & 195 & 55 & 62 & 15 & 7 & 18 & 12 & 11 & 15 \\
143 & 188 & 54 & 52 & 16 & 4 & 16 & 22 & 14 & 10 \\
144* & 205 & 60 & 64 & 17 & 13 & 19 & 10 & 10 & 12 \\
145 & 219 & 58 & 52 & 18 & 15 & 16 & 25 & 18 & 17 \\
146* & 200 & 58 & 60 & 16 & 15 & 21 & 9 & 9 & 12 \\
147 & 212 & 58 & 67 & 16 & 7 & 15 & 22 & 16 & 11 \\
148 & 212 & 63 & 60 & 17 & 15 & 22 & 15 & 7 & 13 \\
149 & 221 & 66 & 65 & 19 & 3 & 13 & 25 & 15 & 15 \\
150* & 204 & 65 & 60 & 17 & 22 & 24 & 5 & 9 & 2 \\
151 & 219 & 67 & 62 & 19 & 9 & 15 & 21 & 14 & 12 \\
152 & 243 & 65 & 64 & 21 & 23 & 25 & 16 & 12 & 17 \\
153 & 263 & 72 & 67 & 25 & 22 & 19 & 8 & 21 & 29 \\
154 & 270 & 80 & 66 & 15 & 15 & 21 & 7 & 40 & 26 \\
155* & 285 & 95 & 80 & 15 & 25 & 25 & 15 & 15 & 15 \\
156* & 300 & 80 & 95 & 33 & 33 & 21 & 21 & 10 & 7 \\
157* & 315 & 90 & 90 & 30 & 30 & 20 & 30 & 15 & 10 \\
158* & 330 & 85 & 130 & 45 & 35 & 7 & 15 & 10 & 3 \\
159* & 345 & 120 & 70 & 60 & 53 & 15 & 7 & 10 & 10 \\
160 & 360 & 50 & 90 & 30 & 70 & 60 & 10 & 10 & 40 \\
161 & 375 & 70 & 85 & 45 & 22 & 25 & 50 & 31 & 47 \\
162* & 390 & 55 & 99 & 70 & 53 & 34 & 30 & 12 & 37 \\
163* & 405 & 130 & 111 & 35 & 81 & 21 & 12 & 7 & 8 \\
\bottomrule
\end{tabular}
}
\scriptsize\caption*{(*) Solver did not produce an initial integer solution}
\end{table}


\begin{table}[!htp]\centering
\caption{Generators parameters from \citet{Morales-Espania2013}.}\label{tab:parametersMorales}
\scriptsize
\scalebox{0.9}{
\begin{tabular}{lrrrrrrr|rrrrr}\toprule
\multicolumn{7}{c}{Technical Information} &\multicolumn{4}{c}{Cost coeficients} \\\cmidrule{2-12}
\multirow{2}{*}{$g$} & $\overline{P}_{g}$ & $\underline{P}_{g}$ & $UT_{g}/DT_{g}$ &$RU_{g}/RD_{g}$ &$U_g/D_g$ &${p}_{g,0}$ &$TC_{g}$ &$C_{g}^{\text{R}}$ &$C{^{l=1}_{g}}^{\dagger}$ &${C^{\text{S}}_{g,s=1}}^{\ddag}$ & ${C^{\text{S}}_{g,s=2}}^{\ddag}$ \\\cmidrule{2-12}
  &[MW] &[MW] &[h] &[MW/h] &[h] &[MW] &[h] &[\$/h] &[MW/h] &[\$] &[\$] \\\cmidrule{1-12}
1 &455  &150  &8   &225    &0   &455  &14  &1000 &16.19 &4500 &9000 \\
2 &455  &150  &8   &225    &0   &245  &14  &970  &17.26 &5000 &10000 \\
3 &130  &20   &5   &50     &0   &0    &10  &700  &16.60 &550  &1100 \\
4 &130  &20   &5   &50     &0   &0    &10  &680  &16.50 &560  &1120 \\
5 &162  &25   &6   &60     &0   &0    &11  &450  &19.70 &900  &1800 \\
6 &80   &20   &3   &60     &0   &0    &8   &370  &22.26 &170  &340 \\
7 &85   &25   &3   &60     &0   &0    &6   &480  &27.74 &260  &520 \\
8 &55   &10   &1   &135    &0   &0    &2   &660  &25.92 &30   &60 \\\midrule
\multicolumn{12}{c}{$^\dagger$ The subindex $s$=1 corresponds to a hot startup, while the subindex $s$=2 corresponds to a cold startup.  } \\
\multicolumn{12}{c}{$^\ddag$ The energy cost is represented as a linear value with a single numerical value.} \\
\multicolumn{12}{c}{The subindex $l$=1 specifically pertains to a step within a staircase function.} \\
\bottomrule
\end{tabular}
}
\end{table}
\end{document}

